\documentclass[aps,prd,twocolumn,superscriptaddress,nofootinbib]{revtex4-2}

\usepackage{amsmath}
\usepackage{mathtools}
\usepackage{amssymb}
\usepackage{amsthm}
\usepackage{booktabs}
\usepackage{tabularx}
\usepackage{array}
\usepackage{graphicx}
\usepackage{xcolor}
\usepackage{mathrsfs}

\newtheorem{theorem}{Theorem}[section]
\newtheorem{proposition}[theorem]{Proposition}
\newtheorem{lemma}[theorem]{Lemma}
\newtheorem{corollary}[theorem]{Corollary}
\newtheorem{conjecture}[theorem]{Conjecture}
\newtheorem{axiom}[theorem]{Axiom}
\theoremstyle{definition}
\newtheorem{definition}[theorem]{Definition}
\theoremstyle{remark}
\newtheorem{remark}[theorem]{Remark}
\newtheorem{claim}[theorem]{Claim}

\DeclareMathOperator{\poly}{poly}
\DeclareMathOperator{\SIZE}{SIZE}
\DeclareMathOperator{\SAT}{SAT}
\newcommand{\NP}{\mathrm{NP}}
\newcommand{\PP}{\mathrm{P}}
\newcommand{\BPP}{\mathrm{BPP}}
\newcommand{\coNP}{\mathrm{coNP}}

\usepackage[hyperfootnotes=false]{hyperref}
\hypersetup{
    pdftitle={Eine entropische Perspektive auf P vs NP: Umformulierung über Kolmogorov-Komplexität},
    pdfauthor={Lukas Geiger},
    pdfsubject={Entropische Reformulierung von P vs NP über ressourcenbeschränkte Kolmogorov-Komplexität},
    pdfkeywords={P vs NP, Kolmogorov-Komplexität, algorithmische Entropie, Zeugenkompression, Berechnungskomplexität},
    colorlinks=true,
    linkcolor=black,
    urlcolor=blue,
    citecolor=black
}
\renewcommand*{\HyperDestNameFilter}[1]{de.#1}

\begin{document}

% ===================================================================
% TITELSEITE
% ===================================================================

\title{Eine entropische Perspektive auf P vs NP: Umformulierung über Kolmogorov-Komplexität}

\author{Lukas Geiger}
\email{https://orcid.org/0009-0005-7296-1534}
\affiliation{Unabhängiger Forscher, Bernau, Deutschland}

\date{5. Mai 2026}

\begin{abstract}
Wir schlagen eine Umformulierung des $\PP$-vs-$\NP$-Problems als \textit{entropische Reformulierung} vor: $\PP \neq \NP$ ist äquivalent zur Nicht-Existenz eines polynomialzeitlichen Zeugen-Selektors, dessen Ausgaben nur logarithmische ressourcenbeschränkte bedingte Komplexität besitzen. Diese Arbeit beweist nicht $\PP \neq \NP$; vielmehr liefert sie eine informationstheoretische Umformulierung, die die präzise Struktur dessen identifiziert, was gezeigt werden muss, und warum klassische Barrierentechniken diesem Ansatz nicht entgegenstehen. Wir führen das Konzept der \textit{algorithmischen Positivität} ein -- eine NP-Sprache ist algorithmisch positiv, wenn sie eine entropisch konvexe, polynomial skalierende Zeugenproduzentenfunktion zulässt -- und charakterisieren die \textit{Entropische Separationsvermutung} (keine NP-vollständige Sprache ist algorithmisch positiv) als \textit{äquivalent} zu $\PP \neq \NP$, womit ESC als Umformulierung des P-vs-NP-Problems in informationstheoretischer Sprache etabliert wird -- nicht als unabhängiges Resultat. Wir führen eine systematische Barrierenanalyse durch und argumentieren, dass Argumente basierend auf unbeschränkter Kolmogorov-Komplexität $K$ den drei klassischen Barrieren (Relativierung, natürliche Beweise, Algebrisierung) nicht natürlich unterliegen, wobei wir anmerken, dass der Barrierenstatus für die in unseren Kernresultaten verwendete ressourcenbeschränkte Variante $K^t$ noch vollständig zu klären ist. Die zentrale strukturelle Einsicht ist die \textit{Zeugen-Entropielücke}: $\PP = \NP$ würde einen kanonischen polynomialzeitlichen Zeugen-Selektor mit $K^t(W(x) \mid x)=O(\log|x|)$ liefern, während $\PP \neq \NP$ impliziert, dass für jede feste polynomiale Zeitschranke und jede feste logarithmische Konstante unendlich viele Instanzen keinen gültigen Zeugen mit derart kurzer ressourcenbeschränkter bedingter Beschreibung besitzen. Die stärkere Nahe-Maximal-Schranke $K^t(w \mid x) \geq |w| - O(\log n)$ wird unten als zusätzliche Brückenhypothese behandelt, nicht als Folgerung aus $\PP \neq \NP$. Wir identifizieren die präzise verbleibende Herausforderung: die Konvertierung dieser Charakterisierung in einen Beweis erfordert eine \textit{Uniformitätsbrücke} -- eine neue Verbindung zwischen algorithmischer Informationstheorie und uniformer Komplexitätstheorie.
\end{abstract}

\keywords{P vs NP, Kolmogorov-Komplexität, algorithmische Entropie, Umformulierung, Barrierenanalyse, Zeugenkompression, Berechnungskomplexität}

\thanks{KI-Werkzeuge (Claude von Anthropic) wurden für mathematische Formalisierung, Literaturverifikation und redaktionelle Überarbeitung eingesetzt. Alle mathematischen Inhalte wurden vom Autor entwickelt, verifiziert und freigegeben.}

\maketitle


% ===================================================================
% 1. EINFUEHRUNG
% ===================================================================
\section{Einführung}
\label{sec:intro}

\subsection{Das Problem}

Die $\PP$-vs-$\NP$-Frage fragt, ob jedes Problem, dessen Lösungen in Polynomialzeit \textit{verifiziert} werden können, auch in Polynomialzeit \textit{gelöst} werden kann. Formal:

\begin{definition}
$\PP$ ist die Klasse der in Polynomialzeit entscheidbaren Sprachen. $\NP$ ist die Klasse der Sprachen, für die Mitgliedschaft einen polynomialzeit-verifizierbaren Zeugen besitzt~\cite{Arora2009}. $\PP = \NP$ genau dann, wenn jede Sprache in $\NP$ auch in $\PP$ liegt. Nach dem Cook-Levin-Theorem~\cite{Cook1971} ist dies äquivalent dazu, dass das Erfüllbarkeitsproblem (SAT) in Polynomialzeit lösbar ist.
\end{definition}

Trotz jahrzehntelanger Bemühungen bleibt das Problem offen. Diese Arbeit argumentiert, dass die Schwierigkeit nicht technischer, sondern \textit{konzeptueller} Natur ist: frühere Ansätze versuchen, untere Schranken über \textit{konstruktive} Methoden zu beweisen (Aufzeigen spezifischer schwerer Instanzen, explizite Schaltkreistrennungen oder algebraische Obstruktionen), die systematisch von drei Barrieren blockiert werden. Wir schlagen eine Alternative vor: $\PP \neq \NP$ über ein \textit{strukturelles Entropieargument} umzuformulieren, das die Unmöglichkeit der Zeugenkompression als fundamentale Obstruktion identifiziert und die präzise verbleibende Lücke benennt, die jeder Beweis schließen muss.

\subsection{Die drei Barrieren}

Drei Meta-Theoreme beschränken den Umfang bekannter Beweistechniken:

\textbf{1. Relativierung} \cite{BakerGillSolovay1975}. Es existieren Orakel $A, B$ mit $\PP^A = \NP^A$ und $\PP^B \neq \NP^B$. Daher muss jeder Beweis von $\PP \neq \NP$ Techniken verwenden, die nicht relativieren -- er muss die interne Struktur der Berechnung ausnutzen, nicht nur das Ein-/Ausgabeverhalten.

\textbf{2. Natürliche Beweise} \cite{RazborovRudich1997}. Wenn starke Pseudozufallsgeneratoren existieren (eine weitverbreitete Annahme), dann kann keine ,,natürliche'' Beweismethode -- eine, die sowohl \textit{konstruktiv} (erkennt schwere Funktionen effizient) als auch \textit{groß} (gilt für eine dichte Menge von Funktionen) ist -- superpolynomiale Schaltkreis-untere-Schranken gegen $\PP/\poly$ beweisen.

\textbf{3. Algebrisierung} \cite{AaronsonWigderson2009}. Selbst nicht-relativierende Techniken, die auf Arithmetisierung basieren (wie in $\mathrm{IP} = \mathrm{PSPACE}$), können $\PP$ nicht von $\NP$ trennen, weil algebraische Erweiterungen von Orakeln existieren, in denen sowohl $\PP = \NP$ als auch $\PP \neq \NP$ gelten.

\subsection{Die Positivitätsalternative}

Wir beobachten eine \textit{oberflächliche strukturelle Parallele} über mehrere große mathematische Probleme hinweg: Durchbrüche gelingen oft nicht durch Konstruktion des gewünschten Objekts, sondern indem gezeigt wird, dass seine Nichtexistenz eine \textit{Positivitätsbedingung} verletzen würde. Die Analogie liegt auf der Ebene der Beweisstrategie (nichtkonstruktive Argumente über nichtnegative Größen), nicht auf der Ebene des mathematischen Inhalts -- die beteiligten Positivitätsstrukturen sind fundamental verschieden:

\begin{center}
\resizebox{\columnwidth}{!}{%
\begin{tabular}{@{}lll@{}}
\hline
Problem & Positive Struktur & Status \\
\hline
Weil-Vermutungen & Frobenius-Eigenwerte & Bewiesen \\
BSD (Rang $\leq 1$) & N\'eron-Tate-Höhe & Bewiesen \\
Hodge-Vermutung & Hodge-Riemann-Form & Offen \\
RH & Spektraler Zeta-Fluss & \cite{GeigerRH2026c} \\
\textbf{P vs NP} & \textbf{Kolmogorov-Komplexität} & \textbf{Diese Arbeit} \\
Breiteres Programm & Freie-Energie-Framework & \cite{FST-Unified2026} \\
\hline
\end{tabular}
}
\end{center}

\noindent\textit{Caveat:} Diese Tabelle illustriert strukturelle Parallelen auf der Ebene der Beweisstrategie; sie behauptet nicht, dass P vs NP durch die Positivitätsperspektive gelöst oder vereinfacht wird. Die beteiligten Positivitätsstrukturen sind in jedem Fall fundamental verschieden.

Die zentrale Einsicht: \textit{Kolmogorov-Komplexität ist die natürliche Positivitätsstruktur für Berechnungskomplexität}. Sie ist nichtnegativ ($K(x) \geq 0$), maschineninvariant (bis auf additive Konstanten) und fundamental nichtkonstruktiv (nicht berechenbar) -- präzise die Eigenschaften, die benötigt werden, um die drei Barrieren für unbeschränktes $K$ zu umgehen. Die ressourcenbeschränkte Variante $K^t$, die für die Kernresultate essentiell ist, erbt einige, aber nicht alle dieser Eigenschaften; ihr Barrierenstatus wird in Abschnitt~\ref{sec:entropy} diskutiert.

Diese Arbeit ist ein \textbf{Reformulierungspaper}: sie etabliert eine Äquivalenz zwischen $\PP \neq \NP$ und einer entropischen Bedingung, beweist aber nicht die Separation selbst. Die Arbeit ist in sich geschlossen; alle Hauptresultate hängen nur von Definitionen und Beweisen innerhalb dieses Dokuments ab. Verbindungen zum breiteren FST-Programm werden als Kontext zitiert, sind aber logisch nicht erforderlich.


% ===================================================================
% 2. BARRIERENANALYSE
% ===================================================================
\section{Systematische Barrierenanalyse}
\label{sec:barriers}

\subsection{Warum konstruktive Methoden scheitern}

Alle drei Barrieren teilen eine gemeinsame Ursache: sie blockieren Beweise, die Berechnungsprobleme als \textit{Black Boxes} behandeln oder die auf \textit{statistischen} Eigenschaften Boolescher Funktionen beruhen.

\begin{itemize}
\item \textbf{Relativierung} blockiert Beweise, die die Berechnung als Ein-/Ausgabeverhalten (Orakelzugriff) behandeln und die interne Struktur ignorieren.
\item \textbf{Natürliche Beweise} blockieren Beweise, die schwere Funktionen über effizient testbare, statistisch dichte Eigenschaften identifizieren.
\item \textbf{Algebrisierung} blockiert Beweise, die nur die polynomiale Struktur arithmetisierter Berechnungen verwenden.
\end{itemize}

Der gemeinsame Faden: alle drei Barrieren gelten für \textit{modellabhängige} Argumente -- Argumente, die innerhalb eines spezifischen Berechnungsmodells (Orakelmaschinen, Schaltkreise, algebraische Erweiterungen) arbeiten, anstatt auf den \textit{intrinsischen Informationsgehalt} der Berechnung zu zielen.

\subsection{Die Spezifikation für einen gültigen Beweis}

Aus der Barrierenanalyse leiten wir notwendige Bedingungen für jeden Beweis von $\PP \neq \NP$ ab:

\begin{enumerate}
\item[\textbf{(S1)}] \textbf{Nicht-relativierend:} Das Argument muss interne Berechnungsstruktur ausnutzen, nicht nur Orakelverhalten.
\item[\textbf{(S2)}] \textbf{Nicht-natürlich:} Das Argument darf keine ,,Größe''-Eigenschaft definieren, die effizient berechenbar ist. Es muss auf einzelne Objekte oder nichtkonstruktive Eigenschaften zielen.
\item[\textbf{(S3)}] \textbf{Nicht-algebrisierend:} Das Argument muss über polynomiale Identitätstests und Arithmetisierung hinausgehen.
\item[\textbf{(S4)}] \textbf{Modellunabhängig:} Das Argument darf nicht von einem spezifischen Berechnungsmodell (Schaltkreise, Turingmaschinen usw.) abhängen, sondern muss auf einer invarianten strukturellen Eigenschaft beruhen.
\item[\textbf{(S5)}] \textbf{Robust:} Das Argument darf nicht zusammenbrechen, wenn man von eingeschränkten Modellen (monotone Schaltkreise, $\mathrm{AC}^0$) zu allgemeiner Berechnung übergeht.
\end{enumerate}

Wir zeigen nun, dass algorithmische Entropie alle fünf Spezifikationen erfüllt.


% ===================================================================
% 3. ALGORITHMISCHE ENTROPIE
% ===================================================================
\section{Algorithmische Entropie als Positivitätsstruktur}
\label{sec:entropy}

\subsection{Kolmogorov-Komplexität}

Sei $U$ eine feste universelle präfixfreie Turingmaschine. Die \textit{Kolmogorov-Komplexität} (algorithmische Entropie) einer endlichen Binärzeichenkette $x$ ist:
\begin{equation}
K(x) := \min\{|p| : U(p) = x\}\,.
\label{eq:kolmogorov}
\end{equation}

Die \textit{bedingte} Kolmogorov-Komplexität von $x$ gegeben $y$ ist:
\begin{equation}
K(x \mid y) := \min\{|p| : U(p, y) = x\}\,.
\label{eq:conditional_K}
\end{equation}

\begin{theorem}[Invarianz {\cite{Kolmogorov1965,Solomonoff1964,Chaitin1966}}]
\label{thm:invariance}
Für je zwei universelle präfixfreie Maschinen $U, U'$:
\begin{equation}
|K_U(x) - K_{U'}(x)| \leq c_{U,U'}\,,
\end{equation}
wobei $c_{U,U'}$ eine von $x$ unabhängige Konstante ist.
\end{theorem}

\begin{theorem}[Unberechenbarkeit]
\label{thm:incomputable}
$K(x)$ ist nicht berechenbar. Genauer: für jede total berechenbare Funktion $f$ existiert $x$ mit $K(x) > f(x)$.
\end{theorem}

\begin{theorem}[Inkompressibilität]
\label{thm:incompressible}
Für jedes $n$ erfüllen mindestens $2^n - 2^{n-c+1} + 1$ Zeichenketten der Länge $n$ die Bedingung $K(x) \geq n - c$.
\end{theorem}

\begin{theorem}[Symmetrie der Information {\cite{LiVitanyi2008}}]
\label{thm:symmetry}
Für alle Zeichenketten $x, y$:
\begin{equation}
K(x, y) = K(x) + K(y \mid x, K(x)) + O(\log(K(x, y)))\,.
\end{equation}
Dies impliziert, dass bedingte Komplexität wohlverhalten ist: die Information im Paar $(x, y)$ zerlegt sich additiv (bis auf logarithmische Terme) in die Information in $x$ plus die zusätzliche Information in $y$ gegeben $x$.
\end{theorem}

\subsection{Warum unbeschränkte Kolmogorov-Komplexität die Barrieren umgeht}

\begin{claim}[Barrierenimmunität]
\label{claim:barrier_immune}
Argumente basierend auf \emph{unbeschränkter} Kolmogorov-Komplexität $K$ erfüllen plausibel die Spezifikationen (S1)--(S5). Wir präsentieren heuristische Argumente für jede; eine vollständig rigorose Behandlung ist eine offene Forschungsfrage (siehe Limitierungen, Abschnitt~\ref{sec:discussion}). \textbf{Wichtiger Vorbehalt:} Die zentralen technischen Resultate dieser Arbeit verwenden \emph{ressourcenbeschränktes} $K^t$, für das die Barrierenimmunität schwächer und teilweise offen ist -- siehe Remark~\ref{rem:barrier-K-vs-Kt} unten.
\end{claim}

\begin{enumerate}
\item[\textbf{(S1)}] \textbf{Nicht-relativierend.} $K(x)$ ist ohne Bezug auf ein Orakel definiert -- es ist eine absolute, feste Größe für jede Zeichenkette $x$. Die Zeugen-Entropielücke (Theorem~\ref{thm:high_entropy}) verwendet $K^t(w \mid x)$ (ressourcenbeschränkte bedingte Komplexität). Wir bemerken eine wichtige Subtilität: das Invarianztheorem für ressourcenbeschränktes $K^t$ ist schwächer als für unbeschränktes $K$ -- der Wechsel der universellen Maschine kann einen \textit{polynomialen} Verlangsamungsfaktor einführen statt nur einer additiven Konstante~\cite{LiVitanyi2008}. Dennoch bleibt das Argument nicht-relativierend, da $K^t$ über ein festes Berechnungsmodell ohne Orakelzugriff definiert ist. Ein relativierender Beweis müsste für alle Orakel $A$ gelten, aber unser Argument zielt auf orakelunabhängige Größen: die Lücke zwischen $O(\log n)$ und $\Omega(|w|)$ ist superpolynomial und somit robust unter polynomialen Verlangsamungsfaktoren.

\item[\textbf{(S2)}] \textbf{Nicht-natürlich.} $K(x)$ ist nicht berechenbar (Theorem~\ref{thm:incomputable}), also kann kein effizientes Verfahren testen, ob eine spezifische Funktion ,,hohe Kolmogorov-Komplexität'' hat. Dies verletzt direkt die ,,Konstruktivitäts''-Anforderung natürlicher Beweise.

\item[\textbf{(S3)}] \textbf{Nicht-algebrisierend.} Algebrisierung betrifft Beweistechniken, die nur die polynomiale Struktur arithmetisierter Berechnungen nutzen. Argumente basierend auf $K^t$ operieren über die volle Stärke der Turingmaschinen-Berechnung, nicht über polynomiale Identitätstests. Ob $K^t$-basierte Argumente formal nicht-algebrisierend im technischen Sinne von~\cite{AaronsonWigderson2009} sind, bleibt eine Forschungsfrage, aber das Entropie-Framework ist nicht natürlich als algebraische Identität ausdrückbar, was nahelegt, dass es dieser Barriere entgeht.

\item[\textbf{(S4)}] \textbf{Modellunabhängig.} Das Invarianztheorem stellt sicher, dass $K(x)$ (bis auf $O(1)$) für alle universellen Maschinen gleich ist -- Turingmaschinen, Schaltkreise, $\lambda$-Kalkül usw.

\item[\textbf{(S5)}] \textbf{Robust.} $K(x)$ hängt nicht vom Berechnungsmodell ab, das zu seiner Definition verwendet wird.
\end{enumerate}

\begin{remark}[Barrierenimmunität: unbeschränktes $K$ vs.\ ressourcenbeschränktes $K^t$]
\label{rem:barrier-K-vs-Kt}
Die oben beanspruchte Barrierenimmunität gilt für unbeschränkte Kolmogorov-Komplexität $K$, die eine absolute, unberechenbare Größe ist, invariant bis auf $O(1)$ unter Wechsel der universellen Maschine. Die zentralen technischen Resultate dieser Arbeit (Theorem~\ref{thm:high_entropy}, die Zeugen-Entropielücke) verwenden jedoch notwendigerweise \emph{ressourcenbeschränktes} $K^t$, das berechenbar ist und dessen Invarianz schwächer ausfällt (polynomiale Verlangsamung statt additiver Konstante). Der Barrierenstatus von $K^t$-basierten Argumenten ist daher weniger eindeutig:
\begin{itemize}
\item \textbf{Relativierung:} $K^t(w \mid x)$ ist ohne Orakel definiert, sodass unsere Argumente nicht-relativierend sind in dem Sinne, dass sie nicht ,,für alle Orakel'' gelten. Jedoch kann $K^{t,A}(w \mid x)$ (mit Orakelzugriff) sich wesentlich von $K^t(w \mid x)$ unterscheiden, und ein vollständiger Barrierenimmunitäts-Beweis für den ressourcenbeschränkten Fall bleibt offen.
\item \textbf{Natürliche Beweise:} $K^t$ ist berechenbar, entgeht also anders als $K$ nicht automatisch der Konstruktivitätsanforderung. Das Argument, dass unser Ansatz natürliche Beweise vermeidet, beruht darauf, dass wir $K^t$ nicht als effizient testbare ,,Größe''-Eigenschaft Boolescher Funktionen verwenden, aber diese Unterscheidung bedarf weiterer Formalisierung.
\item \textbf{Algebrisierung:} Es existiert kein formaler Beweis, dass $K^t$-basierte Argumente im technischen Sinne von~\cite{AaronsonWigderson2009} nicht-algebrisierend sind.
\end{itemize}
Zusammenfassend: Die Barrierenumgehung ist rigoros für unbeschränktes $K$, bleibt aber heuristisch für ressourcenbeschränktes $K^t$. Da die Zeugen-Entropielücke $K^t$ erfordert, ist die Etablierung formaler Barrierenimmunität für den ressourcenbeschränkten Fall ein wichtiges offenes Problem.
\end{remark}


% ===================================================================
% 4. DIE ZEUGEN-ENTROPIELUECKE
% ===================================================================
\section{Die Zeugen-Entropielücke}
\label{sec:gap}

\subsection{NP-Zeugen als komprimierte Information}

Für eine NP-Sprache $L$ hat jede Ja-Instanz $x \in L$ einen polynomiell langen Zeugen $w$, der in Polynomialzeit verifizierbar ist:
\begin{equation}
x \in L \iff \exists w,\, |w| \leq \poly(|x|),\, V(x, w) = 1\,.
\end{equation}

Die bedingte Komplexität des Zeugen gegeben die Instanz misst, wie viel \textit{zusätzliche} Information der Zeuge jenseits der Instanz enthält:
\begin{equation}
\begin{aligned}
K(w \mid x)
&= \text{Zusatzinformation, um } w\\
&\quad \text{aus } x \text{ zu erzeugen.}
\end{aligned}
\end{equation}

\subsection{Was P = NP implizieren würde}

\begin{proposition}[Zeugenkompression unter P = NP]
\label{prop:compression}
Wenn $\PP = \NP$, dann existiert für jede NP-Sprache $L$ und jedes $x \in L$ ein polynomialzeitlicher kanonischer Zeugen-Selektor $A$, der aus $x$ einen gültigen Zeugen $w$ erzeugt:
\begin{equation}
w = A(x)\,.
\end{equation}
Daher gilt für ein Polynom $q$, das von der Laufzeit von $A$ abhängt:
\begin{equation}
K^{q(|x|)}(w \mid x) \leq K(A) + O(\log|x|) = O(\log|x|)\,,
\label{eq:compression}
\end{equation}
da der Algorithmus $A$ ein festes endliches Programm ist, das in Polynomialzeit läuft, und $O(\log|x|)$ Bits die Eingabelänge kodieren.
\end{proposition}

\begin{proof}
Wenn $\PP = \NP$, ist die Suchversion von SAT in Polynomialzeit lösbar (durch Selbstreduzierbarkeit). Gegeben eine erfüllbare Formel $\varphi$, findet der Algorithmus $A$ eine erfüllende Belegung Bit für Bit in Polynomialzeit. Das Programm ,,führe $A$ auf Eingabe $x$ aus'' hat Länge $|A| + O(\log|x|)$, wobei $|A|$ eine Konstante ist.
\end{proof}

\subsection{Die Entropieobstruktion}

\begin{remark}[Unbeschränkte vs.\ ressourcenbeschränkte Komplexität -- warum $K^t$ essentiell ist]
\label{rem:bounded-vs-unbounded}
Die Unterscheidung zwischen $K(w \mid x)$ und $K^t(w \mid x)$ ist entscheidend und muss vor dem Haupttheorem verstanden werden. Für jede entscheidbare Sprache $L$ und jedes $(x, w)$ mit $V(x,w) = 1$ hat die Brute-Force-Suche $K(w \mid x) \leq O(1)$ (das Suchprogramm hat konstante Länge -- es enumeriert einfach alle Kandidaten und prüft jeden). Der informationstheoretische Gehalt von NP-Härte liegt vollständig in der \textit{Zeit}, die zur Extraktion des Zeugen benötigt wird, nicht in der \textit{Beschreibungslänge}. Deshalb ist $K^t$ mit polynomialem $t$ das korrekte Maß: es erfasst die Berechnungskosten der Zeugenextraktion, was die Essenz der $\PP$-vs-$\NP$-Frage ist.
\end{remark}

\begin{theorem}[Superlogarithmische Zeugenobstruktion -- ressourcenbeschränkt, bedingt]
\label{thm:high_entropy}
Wenn $\PP \neq \NP$, dann gilt für SAT: Für jedes Polynom $q$ und jede Konstante $C>0$ existieren unendlich viele erfüllbare Formeln $x$, sodass \textbf{jede} erfüllende Belegung $w$ die Bedingung erfüllt:
\begin{equation}
K^{q(|x|)}(w \mid x) > C \log |x|\,.
\label{eq:high_entropy}
\end{equation}
Das heißt: Für jede feste polynomiale Zeitschranke und jede feste logarithmische Rat-Konstante gibt es unendlich viele Instanzen, deren gültige Zeugen keine so kurze zeitbeschränkte bedingte Beschreibung besitzen. Dieselbe Aussage gilt für jedes NP-vollständige Suchproblem, auf das SAT über eine polynomialzeitliche zeugenerhaltende Suchreduktion reduziert.

Äquivalent (per Kontraposition): Wenn ein Polynom $q$ und eine Konstante $C$ existieren, sodass jede erfüllbare Formel $x$ eine erfüllende Belegung $w$ mit $K^{q(|x|)}(w \mid x) \leq C\log |x|$ besitzt, dann $\PP = \NP$.

Bemerkung zur Quantorenreihenfolge: Das Theorem ist bewusst für jedes feste Paar $(q,C)$ formuliert. Es behauptet \emph{nicht}, dass eine einzige unendliche Instanzmenge allen polynomialen Zeitschranken gleichzeitig widersteht, und es behauptet \emph{nicht} Nahe-Maximalität der Zeugenentropie.
\end{theorem}

\begin{proof}
Wir beweisen die Kontraposition für SAT. Angenommen, ein Polynom $q$ und eine Konstante $C$ haben die Eigenschaft, dass jede erfüllbare Formel $x$ eine erfüllende Belegung $w$ mit $K^{q(|x|)}(w \mid x) \leq C\log |x|$ besitzt. Dann existiert für jede erfüllbare Formel $x$ ein Programm $p_x$ der Länge höchstens $C\log |x|$, das in Zeit $q(|x|)$ eine erfüllende Belegung ausgibt. Es gibt höchstens $|x|^C$ solche Programme. Ein Polynomialzeit-SAT-Suchalgorithmus kann alle Programme dieser Länge enumerieren, jedes für $q(|x|)$ Schritte mit Eingabe $x$ laufen lassen und die Ausgaben verifizieren. Da SAT-Suche durch Selbstreduzierbarkeit polynomialzeitlich äquivalent zur SAT-Entscheidung ist, folgt $\PP=\NP$.

Die Aussage "unendlich viele" folgt aus derselben Kontraposition plus endlichem Hardcoding: Wenn nur endlich viele Formeln die Schranke verletzten, könnte man deren Zeugen in einer endlichen Tabelle speichern und das obige Enumerationsverfahren für alle übrigen Instanzen verwenden. Eine explizite Konstruktion solcher schwerer Instanzen liefert das Argument nicht; das Theorem ist eine nichtkonstruktive Umformulierung.

\textit{Warum unbeschränktes $K$ versagt:} Für jede SAT-Instanz $x$ mit einer erfüllenden Belegung $w$ hat das Brute-Force-Programm ,,enumeriere alle Belegungen, teste jede, gib die erste erfüllende aus'' Länge $O(1)$ und findet $w$ aus $x$ in Zeit $O(2^{|x|})$. Somit $K(w \mid x) \leq O(1)$ -- die unbeschränkte bedingte Komplexität ist trivial klein. Deshalb ist die Ressourcenschranke essentiell: der informationstheoretische Gehalt von NP-Härte liegt in der \textit{Zeit} zur Extraktion des Zeugen, nicht in der Beschreibungslänge.
\end{proof}

\begin{conjecture}[Nahe-maximale Zeugenentropie]
\label{conj:near-maximal-witness-entropy}
Für eine NP-vollständige Suchformulierung und jedes Polynom $q$ existieren unendlich viele Ja-Instanzen $x$, sodass jeder gültige Zeuge $w$ erfüllt:
\begin{equation}
K^{q(|x|)}(w \mid x) \geq |w| - O(\log |x|)\,.
\end{equation}
Diese Nahe-Maximal-Verstärkung würde die scharfe $\Omega(|w|)$-gegen-$O(\log |x|)$-Entropielücke liefern, die in mehreren Brückenzielen heuristisch verwendet wird. Sie folgt nicht aus Theorem~\ref{thm:high_entropy}; sie ist eine zusätzliche offene Hypothese.
\end{conjecture}


\subsection{Die Lücke}

Kombination von Proposition~\ref{prop:compression} und Theorem~\ref{thm:high_entropy}:

\begin{corollary}[Zeugen-Entropielücke -- die zentrale Widerspruchsstruktur]
\label{cor:gap}
Für jedes feste Polynom $q$ und jede logarithmische Konstante $C$ liefert Theorem~\ref{thm:high_entropy} unendlich viele SAT-Instanzen, deren gültige Zeugen in Zeit $q(|x|)$ nicht $C\log |x|$-komprimierbar sind. Für Eindeutig-Zeugen-Instanzen ist die Widerspruchsstruktur am schärfsten:
\begin{itemize}
\item $\PP = \NP$ $\Rightarrow$ ein polynomialzeitlicher kanonischer Zeugen-Selektor $A$ existiert mit $K^{q(|x|)}(A(x) \mid x) \leq O(\log |x|)$ (Proposition~\ref{prop:compression}).
\item $\PP \neq \NP$ $\Rightarrow$ für jedes feste $(q,C)$ existieren unendlich viele Instanzen ohne gültigen Zeugen $w$ mit $K^{q(|x|)}(w \mid x) \leq C\log |x|$ (Theorem~\ref{thm:high_entropy}).
\end{itemize}
Dies ist kein Beweis von $\PP \neq \NP$, sondern eine Charakterisierung der logarithmischen Kompressionsschwelle: $\PP = \NP$ genau dann, wenn ein polynomialzeitlicher kanonischer Zeugen-Selektor mit logarithmischer bedingter Beschreibung für SAT existiert. Die stärkere $\Omega(|w|)$-gegen-$O(\log n)$-Lücke folgt nur unter der Nahe-Maximal-Hypothese (Conjecture~\ref{conj:near-maximal-witness-entropy}).

Die Lücke ist am schärfsten für Instanzen mit einer eindeutigen erfüllenden Belegung; für Instanzen mit mehreren Zeugen schließt das Theorem auf den harten Instanzen jede logarithmisch kurze polynomialzeitliche Zeugenbeschreibung aus.
\end{corollary}

\begin{remark}[Eindeutige vs.\ mehrere Zeugen]
Theorem~\ref{thm:high_entropy} garantiert Instanzen, bei denen \textit{jeder} gültige Zeuge relativ zur festen Zeitschranke und Konstante superlogarithmische bedingte Komplexität hat. Für diese Instanzen gilt die Widerspruchsstruktur unabhängig davon, wie viele Zeugen existieren: jeder Algorithmus $A$ muss \textit{irgendeinen} gültigen Zeugen ausgeben, und keiner der gültigen Zeugen hat die ausgeschlossene logarithmisch kurze Beschreibung.

Für Eindeutig-Zeugen-Instanzen ist die Struktur besonders sauber: $A$ hat keine Wahl, welchen Zeugen er erzeugt. Im Allgemeinen existieren die hochentropischen Instanzen, die durch Theorem~\ref{thm:high_entropy} garantiert werden, jedoch nichtkonstruktiv (per Kontraposition), und es bleibt offen, ob sie explizit konstruiert werden können -- siehe Abschnitt~\ref{sec:obstruction}.
\end{remark}


% ===================================================================
% 5. ALGORITHMISCHE POSITIVITAET
% ===================================================================
\section{Algorithmische Positivität}
\label{sec:positivity}

Wir formalisieren nun das No-Go-Framework.

\subsection{Definition}

\begin{definition}[Algorithmische Positivität]
\label{def:alg_pos}
Eine NP-Sprache $L$ (mit Verifizierer $V$) ist \textit{algorithmisch positiv}, wenn eine \textit{Zeugenproduzenten}-Funktion $W: \{0,1\}^* \to \{0,1\}^*$ existiert, die (AP1) und (AP2) erfüllt:

\begin{enumerate}
\item[\textbf{(AP1)}] \textbf{Polynomiale Skalierung.} $W$ ist in Zeit $\poly(|x|)$ berechenbar, und für jedes $x \in L$: $V(x, W(x)) = 1$.

\item[\textbf{(AP2)}] \textbf{Entropische Konvexität.} Es existiert ein Polynom $q$ so dass für jedes $x \in L$:
\begin{equation}
K^{q(|x|)}(W(x) \mid x) \leq O(\log|x|)\,.
\end{equation}
Die Zeugenproduktion ist ,,entropisch billig'': der Zeuge $W(x)$ fügt höchstens logarithmische ressourcenbeschränkte Information jenseits der Eingabe hinzu. (Bemerkung: Da $W$ ein festes Polynomialzeit-Programm ist, wird das Polynom $q$ durch $W$ bestimmt und hängt nicht von der Instanz $x$ ab.)
\end{enumerate}
\end{definition}

\begin{remark}
Man beachte, dass (AP2) auf den \textit{Zeugen} $W(x)$ zielt, nicht nur auf das Entscheidungsbit. Ein einzelnes Entscheidungsbit $D(x) \in \{0,1\}$ erfüllt trivial $K(D(x) \mid x) \leq 1$, sodass entropische Konvexität für Entscheidungsfunktionen inhaltsleer wäre. Das aussagekräftige Maß ist die ressourcenbeschränkte Komplexität der Produktion eines \textit{Zeugen} -- hier liegt der informationstheoretische Gehalt von NP-Härte.
\end{remark}

\begin{remark}
Jede $\PP$-entscheidbare NP-Sprache ist algorithmisch positiv: wenn $L \in \PP \cap \NP$, dann produziert der Polynomialzeit-Suchalgorithmus (über Selbstreduzierbarkeit für NP-vollständige Probleme, oder direkt für natürliche NP-Probleme) einen Zeugen $W(x)$ mit $K^{q(|x|)}(W(x) \mid x) \leq |W_{\text{prog}}| + O(\log|x|) = O(\log|x|)$.
\end{remark}

\begin{remark}[Heuristische Eigenschaften von P-entscheidbaren NP-Sprachen]
\label{rem:heuristic-ap}
Für NP-Sprachen in $\PP$ erfüllt die Zeugenproduzenten-Funktion $W$ zusätzlich zwei heuristische Eigenschaften, die die ,,strukturelle Einfachheit'' polynomialer Zeugenextraktion erfassen:

\textbf{(H1) Kompositionelle Stabilität.} Für strukturierte Instanzen unter einem wohldefinierten Kompositionsoperator $\circ$ wird der Zeuge für die komponierte Eingabe typischerweise durch die Komponentenzeugen mit höchstens logarithmischem Aufwand bestimmt: $K^{\poly}(W(x_1 \circ x_2) \mid W(x_1), W(x_2)) \leq O(\log(|x_1|+|x_2|))$.

\textbf{(H2) Monotone Stabilität.} Kleine Störungen der Eingabe erzeugen Zeugen beschränkter zusätzlicher Komplexität: $d_H(x,x') \leq 1 \implies K(W(x) \mid W(x'), x) \leq O(\mathrm{poly}(\log|x|))$.

Diese Eigenschaften sind durch die Struktur von $\NP$-vollständigen Problemen motiviert, werden aber im formalen No-Go-Argument (Theorem~\ref{thm:nogo}) nicht verwendet. Sie dienen als zusätzliche heuristische Evidenz für die Entropische Separationsvermutung.
\end{remark}

\subsection{Die Entropische Separationsvermutung}

\begin{conjecture}[Entropische Separation -- ESC]
\label{conj:entropic}
Keine NP-vollständige Sprache ist algorithmisch positiv. Das heißt: für jede NP-vollständige Sprache $L$ und jede Polynomialzeit-Zeugenproduzenten-Funktion $W$ existieren unendlich viele Ja-Instanzen $x \in L$, sodass $W$ keinen gültigen Zeugen für $x$ produziert (d.h. $V(x, W(x)) = 0$), oder $K^{q(|x|)}(W(x) \mid x) \geq \omega(\log|x|)$ für alle Polynome $q$.

Ein stärkeres paralleles Wahrheitstabellen-Brückenziel (siehe Theorem~\ref{thm:slice_nogo}) lautet: Für eine NP-vollständige Sprache $L$ und jedes Polynom $q$ erfüllen unendlich viele $n$ die Bedingung $K^{q(2^n)}(L_n) \geq n$. Diese Wahrheitstabellen-Aussage wird hier nicht als äquivalent zu ESC behauptet; Theorem~\ref{thm:slice_nogo} zeigt, dass sie eine hinreichende Route zu $\PP\neq\NP$ und ein Brückenziel im Stil von Schaltkreisschranken ist.
\end{conjecture}

\begin{proposition}
\label{prop:equiv}
Die Entropische Separationsvermutung~\ref{conj:entropic} impliziert $\PP \neq \NP$.
\end{proposition}

\begin{proof}
Wenn $\PP = \NP$, dann ist für jede NP-vollständige Sprache $L$ die Suchversion in Polynomialzeit lösbar (für SAT über Selbstreduzierbarkeit; für allgemeines NP-vollständiges $L$ durch Reduktion von $L$ auf SAT, Anwendung des SAT-Suchalgorithmus und Rekonstruktion des Zeugen). Dies liefert eine Polynomialzeit-Zeugenproduzenten-Funktion $W$, die (AP1) erfüllt. Da $W$ ein festes Programm in Polynomialzeit ist, gilt $K^{q(|x|)}(W(x) \mid x) \leq |W_{\text{prog}}| + O(\log|x|) = O(\log|x|)$, was (AP2) erfüllt. Daher wären NP-vollständige Sprachen algorithmisch positiv, im Widerspruch zur Vermutung.
\end{proof}

\begin{remark}[ESC ist äquivalent zu $\PP \neq \NP$]
\label{rem:esc-equivalence}
Die Umkehrung von Proposition~\ref{prop:equiv} gilt ebenfalls: $\PP \neq \NP$ impliziert ESC. Äquivalent argumentiert man per Kontraposition: Wenn ESC fehlschlägt, dann besitzt eine NP-vollständige Sprache eine Polynomialzeit-Zeugenproduzenten-Funktion, die (AP1) und (AP2) erfüllt. Das ist bereits ein Polynomialzeit-Suchverfahren für eine NP-vollständige Sprache, und die übliche Such-zu-Entscheidungs-Reduktion liefert $\PP=\NP$. Somit ESC $\Leftrightarrow$ $\PP \neq \NP$: die Entropische Separationsvermutung ist eine \textit{Umformulierung} des P-vs-NP-Problems in informationstheoretischer Sprache, keine unabhängige Verstärkung. Die Wahrheitstabellen-Version unten ist ein stärkeres hinreichendes Brückenziel über Lemma~\ref{lem:p_low_entropy}.
\end{remark}

\subsection{Warum NP-vollständige Probleme algorithmische Positivität widersetzen}

Das Versagen jeder Bedingung für NP-vollständige Probleme:

\textbf{(AP2) versagt:} Für NP-vollständige Probleme erfordert die Produktion eines Zeugen $w$ für $x \in L$ die Extraktion von Information, die nicht uniform logarithmisch relativ zur Instanz komprimierbar ist. Theorem~\ref{thm:high_entropy} zeigt, bedingt auf $\PP \neq \NP$, dass für jede feste polynomiale Zeitschranke und jede logarithmische Konstante unendlich viele SAT-Instanzen keinen gültigen Zeugen innerhalb dieses Beschreibungsbudgets haben. Die stärkere Nahe-Maximal-Schranke $K^t(w \mid x)=\Omega(|w|)$ ist eine zusätzliche Brückenhypothese, kein Theorem. Keine Polynomialzeit-Zeugenproduzenten-Funktion $W$ kann $K^{\poly}(W(x) \mid x) \leq O(\log|x|)$ für alle SAT-Instanzen erreichen.

Als zusätzliche heuristische Evidenz (nicht Teil des formalen Arguments) versagen auch die Eigenschaften (H1)--(H2) aus Remark~\ref{rem:heuristic-ap} für NP-vollständige Probleme:

\textbf{(H1) versagt:} SAT ist nicht kompositionell zerlegbar im geforderten Sinne. Die Erfüllbarkeit einer Konjunktion $\varphi_1 \wedge \varphi_2$ wird nicht durch die Erfüllbarkeit von $\varphi_1$ und $\varphi_2$ einzeln bestimmt (sie teilen Variablen). Diese Nichtzerlegbarkeit ist präzise die Quelle der NP-Härte.

\textbf{(H2) versagt für zufälliges $k$-SAT nahe der Schwelle:} Für zufällige $k$-SAT-Instanzen nahe der Erfüllbarkeitsschwelle zeigt SAT extreme Sensitivität in seiner Suchversion: das Umkehren einer einzigen Klausel kann den gesamten Zeugenraum umstrukturieren. Spezifisch kann nahe dem Phasenübergang (Klausel-zu-Variablen-Verhältnis $\alpha \approx \alpha_c(k)$) eine einzelne Klauselhinzufügung den Lösungscluster in exponentiell viele Komponenten fragmentieren, was Zeugen unbeschränkter zusätzlicher Komplexität erzeugt~\cite{AchlioptasRicci2006,MezardZecchina2002}. \textbf{Geltungsbereich:} Diese Instabilität ist rigoros nur für zufälliges $k$-SAT im Nahe-Schwellen-Regime $\alpha \approx \alpha_c(k)$ etabliert. Für strukturierte, gepflanzte oder unterbeschränkte Instanzen ($\alpha \ll \alpha_c$) kann der Zeugenraum gut verbunden sein, und monotone Stabilität kann gelten. Das heuristische Argument für das Versagen von (H2) als allgemeines NP-Härte-Phänomen beruht daher auf der Annahme, dass das Nahe-Schwellen-Regime repräsentativ für die kombinatorische Schwierigkeit NP-vollständiger Probleme ist.


% ===================================================================
% 6. DAS NO-GO-THEOREM
% ===================================================================
\section{Das No-Go-Theorem}
\label{sec:nogo}

\begin{theorem}[Entropisches No-Go -- Umformulierung]
\label{thm:nogo}
Die folgenden Aussagen sind äquivalent:
\begin{enumerate}
\item[(i)] $\PP \neq \NP$.
\item[(ii)] Die Entropische Separationsvermutung~\ref{conj:entropic} gilt.
\item[(iii)] Keine NP-vollständige Sprache ist in Polynomialzeit entscheidbar. Das heißt: für jede NP-vollständige Sprache $L$ und jeden Polynomialzeit-Algorithmus $A$ existieren Instanzen $x$ mit $A(x) \neq L(x)$.
\end{enumerate}
Die Äquivalenz $(i) \Leftrightarrow (ii)$ zeigt, dass $\PP \neq \NP$ \textit{äquivalent} zu einer Entropieobstruktion ist: Polynomialzeit-Algorithmen können NP-Zeugen nicht unter ihre intrinsische Kolmogorov-Komplexität komprimieren.
\end{theorem}

\begin{proof}
$(ii) \Rightarrow (i)$: Gezeigt in Proposition~\ref{prop:equiv}. Kontraposition: wenn $\PP = \NP$, dann hat jede NP-vollständige Sprache eine Polynomialzeit-Zeugenproduzenten-Funktion (über den SAT-Suchalgorithmus und Polynomialzeit-Reduktionen), die sowohl (AP1) als auch (AP2) erfüllt. Dies macht NP-vollständige Sprachen algorithmisch positiv, im Widerspruch zur ESC.

$(i) \Rightarrow (ii)$: Kontraposition: angenommen, ESC gilt nicht, d.h.\ eine NP-vollständige Sprache $L$ ist algorithmisch positiv. Dann hat $L$ eine Polynomialzeit-Zeugenproduzenten-Funktion $W$ mit $K^{\poly}(W(x) \mid x) \leq O(\log|x|)$. Insbesondere ist $W$ ein Polynomialzeit-Suchalgorithmus für $L$. Da $L$ NP-vollständig ist, folgt $\PP = \NP$.

$(i) \Leftrightarrow (iii)$: (iii) besagt, dass für jede NP-vollständige Sprache $L$ gilt: $L \notin \PP$. Dies ist äquivalent zu $\PP \neq \NP$ aufgrund der Existenz NP-vollständiger Sprachen (wäre ein NP-vollständiges $L$ in $\PP$, dann $\PP = \NP$; umgekehrt würde $\PP = \NP$ jede NP-vollständige Sprache in $\PP$ platzieren).

Der Inhalt des Theorems ist kein neues Trennungsresultat, sondern die \textit{Umformulierung}: $\PP \neq \NP$ ist äquivalent zu einer entropischen Inkompressibilitätsbedingung für NP-Zeugen. Diese Umformulierung identifiziert die präzise Struktur (Zeugenentropie) und die präzise Lücke (Uniformitätsbrücke), die jeder Beweis adressieren muss.
\end{proof}

\subsection{Ausschluss alternativer Szenarien}

\begin{proposition}[Ausgeschlossene Szenarien]
\label{prop:excluded}
Die folgenden Szenarien sind mit der Konjunktion von entropischer Positivität und Kolmogorov-Obstruktion inkompatibel:

\begin{enumerate}
\item[\textbf{(X1)}] \textbf{Polynomialer Algorithmus mit exponentieller Entropieproduktion.} Ein Algorithmus, der während der Berechnung neue Information ,,erfindet''. Ausgeschlossen, weil für jeden deterministischen Algorithmus $A$ mit Laufzeit $t_A(|x|)$ die ressourcenbeschränkte bedingte Komplexität $K^{t_A(|x|)}(A(x) \mid x) \leq |A| + O(\log|x|)$ erfüllt: das Programm ,,führe $A$ auf Eingabe $x$ aus'' hat konstante Beschreibungslänge. Insbesondere beträgt die \textit{bedingte} algorithmische Information von $A(x)$ gegeben $x$ höchstens $O(1)$ Bits auf der Zeitskala von $A$'s eigener Berechnung -- das heißt, $A$ kann keine Information ,,erzeugen'', die nicht bereits implizit in der Eingabe enthalten war. (Für randomisierte Algorithmen liefern die Zufallsbits $r$ zusätzliche Eingabe, sodass $K(A(x,r) \mid x)$ hoch sein kann. Die P-vs-NP-Frage betrifft jedoch deterministische Berechnung; zudem gilt unter weitverbreiteten Derandomisierungsannahmen $\BPP = \PP$~\cite{ImpagliazzoWigderson1997}, sodass Zufall keine zusätzliche Berechnungskraft liefert.)

\item[\textbf{(X2)}] \textbf{Kompositioneller Kollaps von NP.} Alle SAT-Instanzen zerlegen sich in unabhängig lösbare Teile. Ausgeschlossen durch die NP-Vollständigkeit von SAT: wenn SAT sich zerlegte, dann würden alle NP-Probleme sich zerlegen (über Reduktion), aber UNIQUE-SAT-Instanzen mit zufälligen Lösungen zerlegen sich nicht. (Die Relevanz von UNIQUE-SAT wird durch das Valiant-Vazirani-Theorem~\cite{ValiantVazirani1986} gestützt: SAT reduziert auf UNIQUE-SAT unter randomisierten Reduktionen.)

\item[\textbf{(X3)}] \textbf{Glatte Entscheidungslandschaft.} Alle SAT-Instanzen haben monoton stabile Entscheidungsgrenzen. Ausgeschlossen durch Phasenübergangs-Phänomene in zufälligem SAT: nahe der Erfüllbarkeitsschwelle ändert sich die Entscheidung diskontinuierlich unter kleinen Klauselhinzufügungen.
\end{enumerate}
\end{proposition}


% ===================================================================
% 7. RESSOURCENBESCHRAENKTE SCHNITTENTROPIE
% ===================================================================
\section{Ressourcenbeschränkte Schnittentropie}
\label{sec:slices}

Die Zeugen-Entropielücke aus Abschnitt~\ref{sec:gap} zielt auf einzelne Instanzen. Eine sauberere Formulierung zielt auf die \textit{Wahrheitstabelle} der Sprache auf allen Eingaben gegebener Länge.

\subsection{Schnitte einer Sprache}

\begin{definition}[Schnitt]
Für eine Sprache $L \subseteq \{0,1\}^*$ ist der $n$-te \textit{Schnitt} $L_n \in \{0,1\}^{2^n}$ die Wahrheitstabelle von $L$ auf Eingaben der Länge $n$: $(L_n)_i = 1 \iff x_i \in L$, wobei $x_1, \ldots, x_{2^n}$ die $n$-Bit-Zeichenketten in lexikographischer Ordnung sind.
\end{definition}

\begin{definition}[Ressourcenbeschränkte algorithmische Entropie]
Für eine Zeitschranke $t(\cdot)$:
\begin{equation}
\begin{aligned}
K^t(x) := \min\{\, |p| :{}& U(p) = x\\
&\text{in höchstens } t(|x|)\text{ Schritten}\,\}.
\end{aligned}
\label{eq:Kt}
\end{equation}
\end{definition}

\subsection{Das Schnittentropie-Lemma}

\begin{lemma}[P impliziert niedrige Schnittentropie]
\label{lem:p_low_entropy}
Wenn $L \in \PP$, dann existieren ein Polynom $q$ und eine Konstante $c$, sodass für alle $n$:
\begin{equation}
K^{q(2^n)}(L_n) \leq c + \lceil \log n \rceil\,.
\label{eq:p_entropy}
\end{equation}
\end{lemma}

\begin{proof}
Sei $A$ ein Polynomialzeit-Algorithmus, der $L$ in Zeit $n^d$ entscheidet. Das Programm $p$: ,,lese $n$ von der Eingabe; für jedes $x \in \{0,1\}^n$, führe $A(x)$ aus und gib das Ergebnisbit aus'' hat Länge $|A| + O(\log n)$ (um $n$ zu kodieren). Seine Laufzeit ist $2^n \cdot n^d = \poly(2^n)$. Also $K^{q(2^n)}(L_n) \leq |A| + O(\log n) = c + \lceil \log n \rceil$.
\end{proof}

\subsection{Das No-Go-Theorem über Schnittentropie}

\begin{theorem}[Schnittentropie-No-Go]
\label{thm:slice_nogo}
Wenn eine $\NP$-vollständige Sprache $L$ existiert, die die \textit{Entropiehärte-Hypothese} erfüllt:

\textbf{(EH)} Für jedes Polynom $q$ existieren unendlich viele $n$ mit
\begin{equation}
K^{q(2^n)}(L_n) \geq n\,,
\label{eq:eh}
\end{equation}

dann $\PP \neq \NP$.
\end{theorem}

\begin{proof}
Unmittelbar aus Lemma~\ref{lem:p_low_entropy}: $L \in \PP$ würde $K^{\poly(2^n)}(L_n) = O(\log n)$ ergeben, im Widerspruch zu~\eqref{eq:eh} für große $n$.
\end{proof}

\begin{remark}
Die Entropiehärte-Hypothese (EH) ist eine konkrete, wohldefinierte Vermutung über eine spezifische NP-vollständige Sprache (z.B.\ SAT). Sie behauptet, dass die Wahrheitstabelle von SAT auf $n$-Bit-Formeln nicht auf weniger als $n$ Bits komprimiert werden kann, selbst wenn dem Kompressionsalgorithmus $\poly(2^n)$ Zeit erlaubt wird.

Dies ist eine saubere hinreichende No-Go-Formulierung: die ,,leichte Richtung'' ($L \in \PP \Rightarrow$ niedrige Entropie) ist ein Theorem; die ,,schwere Richtung'' (EH für SAT) ist eine stärkere Wahrheitstabellen-/Schaltkreis-Schranken-Route zu $\PP \neq \NP$, aber keine Folgerung aus $\PP \neq \NP$ allein.

Beachte, dass der Schwellwert $n$ in EH viel kleiner als die Wahrheitstabellenlänge $2^n$ ist: ein zufälliger String der Länge $2^n$ hat $K \approx 2^n$. Der Schwellwert $n$ ist gewählt, weil Lemma~\ref{lem:p_low_entropy} für $L \in \PP$ die Schranke $K^{\poly(2^n)}(L_n) = O(\log n)$ liefert; jeder Schwellwert $\omega(\log n)$ würde genügen, um von P zu separieren, und $n$ liefert eine saubere, moderate untere Schranke, die bereits superpolynomial in der Eingabelänge ist.
\end{remark}

\begin{remark}[Kodierungssensitivität von EH]
\label{rem:encoding}
Die Entropiehärte-Hypothese hängt von der Kodierung der SAT-Instanzen als Binärzeichenketten ab. Verschiedene Kodierungen liefern Wahrheitstabellen $\SAT_n$ unterschiedlicher Länge und potentiell unterschiedlicher Kolmogorov-Komplexität. Die Hypothese ist zu verstehen als: es existiert eine \emph{Standard}-Polynomialzeit-Kodierung $\mathrm{enc}: \mathrm{Formeln} \to \{0,1\}^*$, sodass $K^{q(2^n)}(\SAT_n^{\mathrm{enc}}) \geq n$ für unendlich viele $n$. Für ressourcenbeschränktes $K^t$ führt ein Maschinenwechsel zu einem polynomialen Slowdown statt nur einer additiven Konstante. Da EH aber über \emph{alle} Polynome $q$ quantifiziert, beeinflusst ein polynomialer Slowdown in der Zeitschranke die Hypothese nicht: wenn $K^{q(2^n)}(\SAT_n) \geq n$ für alle $q$ gilt, bleibt dies nach Reparametrisierung erhalten. Die \emph{Länge} der Wahrheitstabelle (und damit der Schwellwert $n$) ist kodierungsabhängig.
\end{remark}

\begin{remark}[Verbindung zu Schaltkreisschranken]
\label{rem:circuit-connection}
Die Entropiehärte-Hypothese ist mit Schaltkreisbeschreibungen verwandt, aber der Schwellwert $n$ ist zu schwach, um für sich genommen übliche nichtuniforme untere Schranken zu implizieren. Eine Boolesche Funktion $f_n$, die von einem Schaltkreis der Größe $s$ berechnet wird, hat eine Wahrheitstabelle, die in $O(s \log s)$ Bits beschreibbar ist, also $K^{O(s \cdot 2^n)}(f_n) \leq O(s \log s) + O(\log n)$. Daher:
\begin{itemize}
\item EH mit Schwellwert $n$ schließt nur Schaltkreisbeschreibungen mit weniger als $n$ Bits aus. Quantitativ folgt für Schaltkreise, die als direkte Beschreibung von $\SAT_n$ dienen, $s\log s=\Omega(n)$, also nur $s=\Omega(n/\log n)$ bis auf Konstanten.
\item EH impliziert \emph{nicht} $\NP \not\subseteq \PP/\poly$ oder $\NP \not\subseteq \mathrm{SIZE}(2^{o(n)})$. Solche Schlüsse bräuchten deutlich stärkere Wahrheitstabellen-Inkompressibilität oder ein zusätzliches Theorem, das uniforme SAT-Schnittkompression in nichtuniforme Schaltkreisuntergrenzen übersetzt.
\end{itemize}
So bleibt das Entropie-Framework mit der Schaltkreiskomplexität kompatibel, ohne zu überschätzen, was der moderate EH-Schwellwert leistet.
\end{remark}

\begin{remark}[SAT-Schnittentropie am Phasenübergang]
\label{rem:sat-entropy}
Als rein illustrative Übung (keine Evidenz für oder gegen ESC) wurde die Shannon-Entropie $H_{\mathrm{slice}}$ der Wahrheitstabelle zufälliger 3-SAT-Instanzen am Phasenübergang $\alpha \approx 4{,}267$ für $n = 8, 10, 12, 14, 16$ berechnet. Die Entropie wächst von $H \approx 1{,}5$ ($n=8$) auf $H \approx 2{,}3$ ($n=14$). \textbf{Einschränkungen:} (i) Nur fünf Datenpunkte liegen vor; das ist viel zu wenig, um asymptotisches Wachstum zu bestimmen -- jede monotone Funktion, auch logarithmisches Wachstum, ist mit diesen Daten vereinbar. (ii) Die Shannon-Entropie $H_{\mathrm{slice}}$ der Wahrheitstabelle ist \emph{nicht} dasselbe wie die Kolmogorov-Komplexität $K^t(\mathrm{SAT}_n)$; der Zusammenhang wird hier nicht formal etabliert. (iii) $n \leq 16$ liegt weit unterhalb des asymptotisch relevanten Bereichs. Eine vollständige Validierung würde $n \geq 30$ erfordern, was rechnerisch nicht durchführbar ist. Skript: \texttt{compute\_sat\_entropy.py}.
\end{remark}

\subsection{Warum (EH) wahr sein sollte}

Die Wahrheitstabelle $\SAT_n$ hat Länge $2^n$ und beschreibt die Erfüllbarkeit aller $n$-Bit-Kodierungen Boolescher Formeln. Es gibt $2^{2^n}$ mögliche Wahrheitstabellen dieser Länge, und $\SAT_n$ ist eine \textit{spezifische}, bestimmt durch eine vollständige mathematische Definition ($\SAT$). Dennoch:

\begin{enumerate}
\item \textbf{Pseudozufällige Struktur.} Die Erfüllbarkeitsgrenze im Raum der Formeln zeigt Phasenübergangsverhalten \cite{Monasson1999}: nahe der Schwelle hat das Muster erfüllbarer vs.\ unerfüllbarer Instanzen hohe scheinbare Zufälligkeit.

\item \textbf{Keine bekannte Kompression.} Trotz jahrzehntelanger SAT-Solver-Forschung komprimiert keine Methode $\SAT_n$ auf $o(2^n)$ Bits in $\poly(2^n)$ Zeit. Die besten Algorithmen (CDCL-Solver) laufen im schlimmsten Fall in Zeit $2^{\Theta(n)}$.

\item \textbf{Hardness-vs-Randomness-Evidenz.} Explizite Wahrheitstabellen-Inkompressibilität ist eng mit Schaltkreisschranken und Derandomisierung verbunden~\cite{ImpagliazzoWigderson1997}. EH ist daher eine natürliche Verstärkung im Umfeld von PRG- und MCSP-Routen; ihr Scheitern würde aber nur diese konkrete Wahrheitstabellen-Route schließen, nicht $\PP\neq\NP$ widerlegen.
\end{enumerate}


% ===================================================================
% 8. DIE PRAEZISE OBSTRUKTION
% ===================================================================
\section{Die präzise Obstruktion}
\label{sec:obstruction}

\subsection{Die Uniformitätslücke}

Das Hauptargument (Abschnitt~\ref{sec:gap}) hat eine präzise Lücke, die wir nun identifizieren.

\textbf{Was bewiesen ist:} Es \textit{existieren} NP-Instanzen, deren Zeugen hohe bedingte Kolmogorov-Komplexität haben (Theorem~\ref{thm:high_entropy}). Dies ist eine \textit{nichtuniforme} Aussage über einzelne Zeichenketten.

\textbf{Was benötigt wird:} Eine \textit{uniforme} Aussage: für jeden Polynomialzeit-Algorithmus $A$ existieren unendlich viele Instanzen $x$, bei denen $A(x)$ versagt. Dies erfordert die Konvertierung der existenziellen Entropieschranke in eine \textit{adversarische} Schranke gegen alle effizienten Algorithmen gleichzeitig.

\begin{definition}[Die Uniformitätsbrücke]
\label{def:uniformity}
Sei $L$ eine NP-vollständige Sprache mit Verifizierer $V$. Die \textit{Uniformitätsbrücke} ist die Aussage: für jede in Polynomialzeit berechenbare Funktion $A: \{0,1\}^* \to \{0,1\}^*$ existieren unendlich viele Ja-Instanzen $x \in L$, sodass $A$ keinen gültigen Zeugen erzeugt:
\begin{equation}
\forall A \in \mathrm{FP},\, \exists^\infty x \in L:\, V(x, A(x)) = 0\,.
\end{equation}
Äquivalent in entropischen Termen: für jedes Polynomialzeit-$A$ und jedes Polynom $q$ existieren unendlich viele $x \in L$, sodass jeder gültige Zeuge $w$ für $x$ die Bedingung $K^{q(|x|)}(w \mid x) > |A| + O(\log|x|)$ erfüllt.

Beachte: die Uniformitätsbrücke ist als Suchaussage formuliert ($A$ muss einen \textit{Zeugen} erzeugen, d.h.\ $A \in \mathrm{FP}$ und $V(x, A(x)) = 1$). Die Verbindung zum Entscheidungsproblem $\PP \neq \NP$ folgt aus der NP-Vollständigkeit von $L$: für NP-vollständige Sprachen sind Such- und Entscheidungsversion polynomialzeit-äquivalent (über Selbstreduzierbarkeit für SAT, oder allgemeiner über Cook-Levin plus Polynomialzeit-Reduktionen).
\end{definition}

Die Uniformitätsbrücke ist das P-vs-NP-Analogon von:
\begin{itemize}
\item der ,,höheren Gross-Zagier-Formel'' in BSD (Kopplung analytischen Rangs zu algebraischem Rang),
\item der ,,schweren Richtung'' in der Hodge-Vermutung (Nachweis, dass Positivität der Hodge-Riemann-Form Algebraizität impliziert).
\end{itemize}

In jedem Fall ist die Positivitätsstruktur etabliert, aber ein \textit{Kopplungstheorem} wird benötigt, um Positivität in ein definitives Resultat umzuwandeln.

\subsection{Bekannte Teilresultate}

Die Uniformitätsbrücke ist in eingeschränkten Situationen teilweise etabliert:

\begin{enumerate}
\item \textbf{Resolution-Beweissysteme:} Exponentielle untere Schranken für Resolution-Widerlegungen zufälliger $k$-SAT-Instanzen nahe der Schwelle \cite{BenSasson2001}. Dies beweist, dass resolutionsbasierte SAT-Solver diese Instanzen nicht effizient lösen können.

\item \textbf{Schaltkreise beschränkter Tiefe:} $\SAT \notin \mathrm{AC}^0$ (Furst-Saxe-Sipser 1984, Ajtai 1983; siehe~\cite{Hastad1987} für optimale Schranken). Dies beweist die Entropieobstruktion für Schaltkreise konstanter Tiefe.

\item \textbf{Monotone Schaltkreise:} Exponentielle untere Schranken für monotone Schaltkreise, die CLIQUE berechnen \cite{Razborov1985}; verwandte Schranken beschränkter Tiefe \cite{Razborov1987}.

\item \textbf{Algebraische Beweissysteme:} Exponentielle untere Schranken für Nullstellensatz- und Polynomialkalkül-Widerlegungen \cite{Grigoriev2001}.
\end{enumerate}

Jeder dieser Fälle ist eine \textit{modellspezifische} Instantiierung der Entropieobstruktion. Die Herausforderung besteht darin, sie für \textit{allgemeine} Polynomialzeit-Berechnung zu beweisen -- was präzise das $\PP \neq \NP$-Problem ist.

\subsection{Warum die Lücke schwer ist}

Die Uniformitätslücke ist schwer aus demselben Grund, aus dem $\PP$ vs $\NP$ schwer ist: sie erfordert eine Aussage über \textit{alle} Polynomialzeit-Algorithmen. Im uniformen Modell ist dies ein universeller Quantor über eine abzählbar unendliche, rekursiv aufzählbare Menge (alle Turingmaschinen mit polynomialer Laufzeit). Im nichtuniformen Modell ($\PP/\poly$) erstreckt sich der Quantor über alle Schaltkreisfamilien polynomialer Größe, eine überabzählbare Menge.

Der Entropieansatz \textit{reduziert} das Problem, \textit{löst} es aber nicht:
\begin{quote}
\textit{P vs NP ist äquivalent zur Frage, ob Kolmogorov-Komplexität für NP-Zeugenstrukturen polynomial komprimiert werden kann. Das Entropie-Framework identifiziert, was gezeigt werden muss (Zeugen-Inkompressibilität) und warum es wahr sein sollte (zufällige Zeugen existieren), aber die Konvertierung in einen Beweis gegen alle Algorithmen bleibt die zentrale Herausforderung.}
\end{quote}

\subsection{Vier Kandidatenbrücken-Strategien}

Wir identifizieren vier konkrete Forschungsrichtungen zum Schließen der Uniformitätslücke:

\textbf{(B1) Dichte Familien mit eindeutigen zufälligen Zeugen.} Konstruiere \textit{dichte} Familien $S_n \subseteq \{0,1\}^{\poly(n)}$ mit $|S_n|/2^{\poly(n)}$ nicht verschwindend, sodass jedes $x \in S_n$ einen eindeutigen gültigen Zeugen $w_x$ mit $K^{\poly(n)}(w_x \mid x) \geq |w_x| - O(\log n)$ hat. Dichte stellt sicher, dass kein Polynomialzeit-Algorithmus alle schweren Instanzen ,,vermeiden'' kann. Die Konstruktion erfordert effiziente Reduktionen von Kolmogorov-zufälligen Zeichenketten auf SAT-Instanzen mit eindeutigen Lösungen -- eine nichttriviale, aber plausible codierungstheoretische Aufgabe.

\textbf{(B2) Selbstreferenzielle Diagonalisierung.} Konstruiere für jede Polynomialzeit-Maschine $A$ eine Instanz $x_A$, deren eindeutiger gültiger Zeuge $w$ eine Beschreibung von $A$ selbst enthält (oder kodiert), sodass $K^{\poly(|x_A|)}(w \mid x_A) \geq |A|$. Dies erzeugt eine selbstreferenzielle Obstruktion: $A$ kann keinen Zeugen erzeugen, der ,,über sich selbst'' ist, ohne seine eigene Beschreibungslänge zu überschreiten. Der Ansatz kombiniert Gödel-artige Diagonalisierung mit Kolmogorov-Komplexität und vermeidet Relativierung, weil die Selbstreferenz auf die \textit{Beschreibung} der Maschine zielt, nicht auf ihr Orakelverhalten.

\textbf{(B3) Ressourcenbeschränkte Kolmogorov-untere-Schranken für Schnitte.} Beweise, dass für SAT (oder eine andere NP-vollständige Sprache) die ressourcenbeschränkte Komplexität $K^{\poly(2^n)}(\SAT_n) \geq 2^{n-o(n)}$ für unendlich viele $n$ erfüllt. Dies ist eine \textit{uniforme} Aussage über die Wahrheitstabelle und impliziert direkt $\PP \neq \NP$ über das Schnittentropie-No-Go (Theorem~\ref{thm:slice_nogo}). Die Strategie verbindet sich mit Beweiskomplexität: wenn jeder kurze Beweis der Erfüllbarkeit einer Instanz hohe Kolmogorov-Komplexität hat, dann kann kein Polynomialzeit-Algorithmus systematisch solche Beweise erzeugen.

\textbf{(B4) Das Minimum Circuit Size Problem (MCSP).} Jüngere Arbeiten verbinden ressourcenbeschränkte Kolmogorov-Komplexität mit Schaltkreiskomplexität über das \textit{Minimum Circuit Size Problem} ($\mathrm{MCSP}$): gegeben eine Wahrheitstabelle $t \in \{0,1\}^{2^n}$ und einen Parameter $s$, entscheide, ob ein Schaltkreis der Größe $\leq s$ existiert, der $t$ berechnet. Allender et al.~\cite{Allender2006} etablierten tiefe Verbindungen zwischen $\mathrm{MCSP}$ und $K^t$. Während $\mathrm{MCSP} \in \NP$, ist sein Komplexitätsstatus offen: es ist nicht bekannt, ob es $\NP$-vollständig ist, aber der Beweis $\mathrm{MCSP} \in \PP$ hätte starke Derandomisierungskonsequenzen. Starke MCSP-Untergrenzen oder Härteresultate können in Schaltkreisuntergrenzen-Programme einspeisen; sie implizieren die Entropie-Härtehypothese erst nach einem zusätzlichen Schritt, der hinreichend starke SAT-Schnitt-Wahrheitstabellen-Inkompressibilität liefert. Die Kette ist daher konditional statt automatisch: \textit{MCSP-Härte $\Rightarrow$ Schaltkreisuntergrenzen $\Rightarrow$ Wahrheitstabellen-Inkompressibilität $\Rightarrow$ EH}, mit quantitativen Verlusten an jedem Pfeil.


% ===================================================================
% 9. BEWEISKOMPLEXITAET ALS TESTFELD
% ===================================================================
\section{Beweiskomplexität als Testfeld}
\label{sec:proof_complexity}

\subsection{SAT-Algorithmen und kurze Beweise}

Eine tiefe Verbindung besteht zwischen Algorithmen und Beweisen:
\begin{equation}
\begin{aligned}
\SAT \in \PP
\implies{}& \text{jede unerfüllbare Formel}\\
& \text{hat eine kurze Widerlegung,}
\end{aligned}
\end{equation}
in einem Beweissystem, das mindestens so stark wie Extended Frege ist.

\begin{proposition}[Beweiskomplexitätsformulierung -- Cook {\cite{Cook1975}}]
\label{prop:proof_complexity}
Wenn $\PP \neq \NP$, dann existiert kein polynomial beschränktes Beweissystem für Tautologien. Äquivalent: $\PP = \NP$ genau dann, wenn ein Beweissystem existiert, in dem jede Tautologie der Länge $n$ einen Beweis der Länge $\poly(n)$ hat.
\end{proposition}

Dies ist Cooks Umformulierung~\cite{Cook1975}: der Beweis superpolynomialer unterer Schranken für die Beweislänge in jedem propositionslogischen Beweissystem ist äquivalent zum Nachweis $\PP \neq \coNP$ (was aus $\PP \neq \NP$ folgt).

\subsection{Entropie von Beweisen}

Das Entropie-Framework verbindet sich mit Beweiskomplexität über die folgende Frage: für eine zufällige unerfüllbare Formel $\varphi$ und ihre kürzeste Widerlegung $\pi$ in einem gegebenen Beweissystem, wie groß ist $K(\pi \mid \varphi)$? Wenn $\pi$ ,,algorithmisch zufällig'' relativ zu $\varphi$ ist -- d.h. $K(\pi \mid \varphi)$ nahe an $|\pi|$ liegt -- dann kann kein effizienter Algorithmus systematisch solche Widerlegungen erzeugen.

Dies liefert ein konkretes Testfeld: wenn man zeigen kann, dass für zufälliges unerfüllbares $3$-SAT nahe der Schwelle die kürzeste Widerlegung in Extended Frege sowohl superpolynomiale Länge als auch hohe bedingte Kolmogorov-Komplexität hat, wäre dies starke Evidenz für die Entropieobstruktion.

\begin{conjecture}[Resolutionsbreite und Kolmogorov-Komplexität]
\label{conj:resolution-width}
Für unerfüllbare CNF-Formeln $\varphi_n$ auf $n$ Variablen erfüllt die minimale Resolutionsbreite $w(\varphi_n \vdash \bot)$
\[
w(\varphi_n \vdash \bot) \geq \Omega\left( K(\pi^* | \varphi_n) / \log n \right),
\]
wobei $\pi^*$ die kürzeste Resolutionswiderlegung von $\varphi_n$ ist. Insbesondere erfordern Formeln, deren Widerlegungen hohe Kolmogorov-Komplexität haben, breite Klauseln.
\end{conjecture}

\textit{Motivation.} Ben-Sasson und Wigderson~\cite{BenSasson2001} zeigten, dass untere Schranken für die Resolutionsbreite untere Schranken für die Resolutionslänge über $|\pi| \geq 2^{w - O(n)}$ implizieren. Unsere Vermutung kehrt die Richtung um: hohe Komplexität des Beweisobjekts (gemessen durch $K$) erzwingt Breite. Falls wahr, würde dies eine informationstheoretische Charakterisierung der Resolutionsbreite liefern und Beweiskomplexität mit dem ESC-Framework verbinden.

\begin{remark}[Resolutionsbreite als thermodynamische Observable]\label{rem:width-observable}
Die Resolutionsbreite $w(\varphi_n \vdash \bot)$ ist robust unter Zeitbeschränkung: Anders als die Kolmogorov-Komplexität~$K$ ist die Breite einer Resolutionswiderlegung eine \emph{syntaktische} Invariante, die nicht von Rechenressourcen abhängt. Dies macht sie zu einer natürlichen "`thermodynamischen Observablen"' in der Beweiskomplexitäts-Analogie: Große Breite erzwingt, dass jede zeitbeschränkte Beschreibung komplex ist, und liefert damit einen konkreten Weg von $K$-unteren Schranken zu $K^t$-unteren Schranken für spezifische Formelfamilien. Die conditionalen MCSP-Härteresultate von Hirahara und Ilango~\cite{Ilango2025} motivieren MCSP als Brückenobjekt, ohne allgemeine Entscheidbarkeit von~$K$ vorauszusetzen.
\end{remark}

\subsection{\texorpdfstring{Entropiebeweis für $\SAT \notin \mathrm{AC}^0$}{Entropiebeweis für SAT nicht in AC0}}
\label{sec:ac0}

Als konkrete Instantiierung des Entropie-Frameworks geben wir ein entropiebasiertes Argument für das klassische Resultat $\mathrm{PAR} \notin \mathrm{AC}^0$ (Furst-Saxe-Sipser 1984, Ajtai 1983). Da $\mathrm{AC}^0$ unter $\mathrm{AC}^0$-Reduktionen abgeschlossen ist und Parität $\mathrm{AC}^0$-reduzierbar auf SAT ist (über eine Standardreduktion, die die Paritätsbedingung als CNF-Formel mittels konstanttiefer Schaltkreise kodiert), impliziert dies auch $\SAT \notin \mathrm{AC}^0$.

\begin{proposition}[Entropieobstruktion für $\mathrm{AC}^0$]
\label{prop:ac0}
Sei $\{C_n\}$ eine Schaltkreisfamilie konstanter Tiefe mit polynomialer Größe $s(n) = n^{O(1)}$, die eine Boolesche Funktion $f_n: \{0,1\}^n \to \{0,1\}$ berechnet. Dann erfüllt die Wahrheitstabelle $f_n \in \{0,1\}^{2^n}$
\begin{equation}
K^{s(n) \cdot 2^n}(f_n) \leq O(n^c \log n)
\end{equation}
für eine Konstante $c$, die von der Tiefe und Größenschranke abhängt. Im Gegensatz dazu erfüllt eine \textit{zufällige} Boolesche Funktion $g_n$ mit hoher Wahrscheinlichkeit $K(g_n) \geq 2^n - O(1)$. Die Entropielücke zwischen $\mathrm{AC}^0$-berechenbaren Funktionen und zufälligen Funktionen ist somit exponentiell: $O(n^c \log n)$ vs.\ $2^n$.
\end{proposition}

\begin{proof}[Beweisskizze]
Ein Schaltkreis der Tiefe $d$ und Größe $s$ auf $n$ Eingaben wird durch $O(s \cdot \log s)$ Bits beschrieben (Gattertypen und Verdrahtung). Gegeben diese Beschreibung wird die Wahrheitstabelle auf allen $2^n$ Eingaben in Zeit $O(s \cdot 2^n)$ durch Brute-Force-Auswertung berechnet. Somit $K^{O(s \cdot 2^n)}(f_n) \leq O(s \cdot \log s) + O(\log n)$. Für polynomiale Größe $s = n^c$: $K^{n^c \cdot 2^n}(f_n) \leq O(n^c \log n)$.

Das klassische Resultat $\mathrm{PAR}_n \notin \mathrm{AC}^0$ (Furst-Saxe-Sipser 1984, Ajtai 1983, H\r{a}stad 1987) zeigt, dass die Paritätsfunktion nicht durch Schaltkreise polynomialer Größe und konstanter Tiefe berechnet werden kann. Im Entropie-Framework bedeutet dies: jeder Schaltkreis, der $\mathrm{PAR}_n$ berechnet, erfordert superpolynomiale Größe $s(n) = 2^{n^{\Omega(1/d)}}$, sodass $K^{\mathrm{poly}(2^n)}(\mathrm{PAR}_n)$ \textit{nicht} durch $O(n^c \log n)$ beschränkt ist, wenn über $\mathrm{AC}^0$-Schaltkreise berechnet. Die Entropiemethode \textit{beweist} $\mathrm{PAR} \notin \mathrm{AC}^0$ nicht erneut, aber sie \textit{reinterpretiert} die Separation: $\mathrm{AC}^0$-berechenbare Wahrheitstabellen nehmen einen verschwindenden Anteil am Raum aller Booleschen Funktionen ein, und Parität liegt außerhalb dieser niedrig-entropischen Region.
\end{proof}

Dies demonstriert die Entropieperspektive auf eine bekannte Separation. Für allgemeines $\PP/\mathrm{poly}$ würde die analoge Schranke $K^{\mathrm{poly}(2^n)}(f_n) \leq O(n^c \log n)$ für jedes $f \in \PP/\mathrm{poly}$ liefern. Der Beweis, dass $\SAT_n$ diese Schranke übersteigt, würde $\NP \not\subseteq \PP/\mathrm{poly}$ etablieren -- ein großes offenes Problem, das $\PP \neq \NP$ impliziert.

\begin{remark}[H\r{a}stads Schranke als Entropiekapazität]
\label{rem:hastad-entropy}
H\r{a}stads Switching-Lemma~\cite{Hastad1987} kann als \emph{informationstheoretische Kapazitätsschranke} interpretiert werden: nach Anwendung einer zufälligen Restriktion, die jede Variable unabhängig mit Wahrscheinlichkeit $1-p$ fixiert, vereinfacht sich ein $\mathrm{AC}^0$-Schaltkreis der Tiefe $d$ auf den überlebenden Variablen zu einem Entscheidungsbaum der Tiefe $O(\log n)^{d-1}$. In Entropieterminologie: jede Schicht eines $\mathrm{AC}^0$-Schaltkreises kann höchstens $O(\log n)$ Bits Information von ihren Eingaben ,,aggregieren'' (über die Fan-In-Schranke). Nach $d$ Schichten ist die gesamte Informationskapazität $O((\log n)^d)$ Bits -- weit unter den $\Omega(n)$ Bits, die zur Berechnung der Parität benötigt werden. Diese Kapazitätsschranke $O((\log n)^d) \ll n$ ist der entropietheoretische Kern der $\mathrm{PAR} \notin \mathrm{AC}^0$-Separation.
\end{remark}

\subsection{Geometric Complexity Theory}

Mulmuleys Geometric Complexity Theory (GCT) \cite{Mulmuley2011} geht $\PP$ vs $\NP$ über algebraische Geometrie und Darstellungstheorie an. Das GCT-Programm sucht \textit{Obstruktionen} -- darstellungstheoretische Invarianten, die leichte von schweren Funktionen unterscheiden.

Im Entropie-Framework sind GCT-Obstruktionen ein spezifischer Typ \textit{algebraischer Entropiezertifikate}: sie zertifizieren, dass eine Funktion hohe ,,algebraische Komplexität'' hat, indem sie eine Invariante aufzeigen, die polynomiell große Schaltkreise nicht erreichen können.

Die Verbindung zur algorithmischen Positivität: GCT-Obstruktionen sind \textit{nicht-natürlich} (sie sind keine effizient berechenbaren Eigenschaften von Wahrheitstabellen) und \textit{nicht-relativierend} (sie hängen von algebraischer Struktur ab). Sie erfüllen somit die Spezifikationen (S1)--(S2). Ob sie (S3) (Nicht-Algebrisierung) erfüllen, ist eine aktive Forschungsfrage.


% ===================================================================
% 10. VERWANDTE ARBEITEN
% ===================================================================
\section{Verwandte Arbeiten}
\label{sec:related}

Die Verbindung zwischen Kolmogorov-Komplexität und Berechnungskomplexität wurde ausführlich untersucht. Wir ordnen unsere Arbeit kurz in diese Literatur ein.

\textbf{Ressourcenbeschränkte Kolmogorov-Komplexität und Schaltkreisschranken.}
Allender et al.~\cite{Allender2006} etablierten tiefe Verbindungen zwischen dem Minimum Circuit Size Problem (MCSP), ressourcenbeschränkter Kolmogorov-Komplexität $K^t$ und Derandomisierung. Sie zeigten, dass $K^t$ signifikante Information über Schaltkreiskomplexität kodieren kann: eine Wahrheitstabelle $f_n$, die von Schaltkreisen der Größe $s$ berechnet wird, erfüllt $K^{O(s \cdot 2^n)}(f_n) \leq O(s \log s)$ (die Zeitschranke $O(s \cdot 2^n)$ berücksichtigt die Auswertung des Schaltkreises auf allen $2^n$ Eingaben), was Slice-Entropie direkt mit Schaltkreisschranken verbindet. Allender~\cite{Allender2001} überblickte, wie Kolmogorov-Komplexität und Derandomisierung interagieren, und argumentierte, dass das Verständnis der Komplexität ,,zufälliger Zeichenketten'' (Strings mit hohem $K^t$) zentral für die Lösung von P vs NP ist.

\textbf{Gleichförmige Untergrenzen über probabilistisches $K^t$.}
Santhanam~\cite{Santhanam2023} schlug einen algorithmischen Ansatz zu gleichförmigen Untergrenzen mittels probabilistischer zeitbeschränkter Kolmogorov-Komplexität ($\mathrm{pK}^{\mathrm{poly}}$) vor. Sein Ansatz liefert \emph{tatsächlich neue Resultate} (gleichförmige $\mathrm{AC}^0$-Untere Schranken) aus nichttrivialen Sampling-Algorithmen, ohne sich auf Hierarchiesätze zu stützen. Die Uniformitätsbrücke unseres Papers adressiert dieselbe konzeptuelle Lücke, die Santhanam mit konkreten algorithmischen Werkzeugen bearbeitet; seine Arbeit zeigt, dass diese Lücke in eingeschränkten Settings teilweise schließbar ist.

\textbf{MCSP und NP-Härte.}
Hirahara~\cite{Hirahara2018} zeigte, dass die NP-Härte von MCSP unter bestimmten Reduktionen neue Schaltkreisschranken implizieren würde, was einen konkreten Pfad von $K^t$-basierter Härte zu P-vs-NP-Trennungen liefert. Die Kette ,,Härte von MCSP $\Rightarrow$ Schaltkreisschranken $\Rightarrow$ Wahrheitstabellen-Inkompressibilität $\Rightarrow$ EH'' motiviert unsere Brückenstrategie B4.

\textbf{Natürliche Beweise und Pseudozufälligkeit.}
Die Natural-Proofs-Barriere~\cite{RazborovRudich1997} ist eng mit der Unberechenbarkeit von $K$ verbunden: eine ,,natürliche'' Eigenschaft muss effizient berechenbar sein, aber hohe Kolmogorov-Komplexität ist es nicht. Unser Ansatz nutzt diese Verbindung systematisch. Impagliazzo und Wigderson~\cite{ImpagliazzoWigderson1997} zeigten, dass starke Schaltkreisschranken Derandomisierung implizieren ($\mathrm{P} = \mathrm{BPP}$), was in Impagliazzos ,,Fünf-Welten''-Framework~\cite{Impagliazzo1995} Heuristica oder darüber entspricht. EH ist daher als starkes Wahrheitstabellen-Brückenziel im Hardness-vs-Randomness-Umfeld zu lesen; ein Scheitern von EH würde uns nicht schon in Algorithmica platzieren.

\textbf{Abgrenzung von früheren Arbeiten.}
Die Hauptneuheit dieser Arbeit ist nicht, neue Verbindungen zwischen Kolmogorov-Komplexität und P vs NP herzustellen -- solche Verbindungen sind seit Orponen et al.~\cite{OrponenKoSchoeningWatanabe1994} und Fortnow~\cite{Fortnow2000} bekannt und wurden von Allender~\cite{Allender2001,Allender2006}, Hirahara~\cite{Hirahara2018}, Liu--Pass~\cite{LiuPass2020} und Santhanam~\cite{Santhanam2023} weiterentwickelt. Unser Beitrag ist die spezifische \emph{Darstellung}: die Rahmung als No-Go-Theorem über algorithmische Positivität, die systematische Barrierenanalyse (für $K$ streng, für $K^t$ heuristisch), die Identifikation der Uniformitätsbrücke als präzise verbleibende Lücke und die Verbindung zur Beweiskomplexität über die Resolution-Width-Vermutung.

% ===================================================================
% 11. DISKUSSION
% ===================================================================
\section{Diskussion}
\label{sec:discussion}

\subsection{Was wurde erreicht}

\begin{enumerate}
\item \textbf{Umformulierung.} $\PP$ vs $\NP$ wird als entropisches No-Go-Theorem umformuliert: Polynomialzeit-Algorithmen können NP-Zeugen nicht unter ihre Kolmogorov-Komplexität komprimieren.

\item \textbf{Barrierenanalyse.} Argumente basierend auf \emph{unbeschränkter} Kolmogorov-Komplexität $K$ erfüllen alle fünf Barrierenspezifikationen: nicht-relativierend, nicht-natürlich, nicht-algebrisierend, modellunabhängig und robust. Für die in den Kernresultaten verwendete \emph{ressourcenbeschränkte} Variante $K^t$ ist der Barrierenstatus heuristisch: $K^t$ ist berechenbar und entgeht daher nicht automatisch der Natural-Proofs-Barriere; formale Nicht-Algebrisierung bleibt offen (siehe Remark~\ref{rem:barrier-K-vs-Kt}).

\item \textbf{Zeugen-Entropielücke.} Bedingt auf $\PP \neq \NP$ existieren SAT-Instanzen, deren gültige Zeugen innerhalb jeder festen polynomialen Zeitschranke und logarithmischen Konstante nicht logarithmisch komprimierbar sind. Die stärkere Behauptung nahe-maximaler bedingter Kolmogorov-Komplexität ist als Conjecture~\ref{conj:near-maximal-witness-entropy} isoliert.

\item \textbf{Algorithmische Positivität.} Das Konzept identifiziert zwei zentrale strukturelle Eigenschaften (Polynomialzeit-Zeugenextraktion und entropische Billigkeit der Zeugen), die $\PP$-entscheidbare NP-Sprachen erfüllen und NP-vollständige Probleme verletzen, ergänzt durch heuristische Eigenschaften (kompositionelle und monotone Stabilität der Zeugen), die zusätzliche Separationsevidenz liefern.
\end{enumerate}

\subsection{Was offen bleibt}

Das zentrale offene Problem -- die \textit{Uniformitätsbrücke} -- ist die Konvertierung existenzieller Entropieschranken (,,es existieren schwere Instanzen'') in universelle Algorithmenschranken (,,kein Algorithmus löst alle Instanzen''). Dies ist die P-vs-NP-spezifische Instanz des allgemeinen Musters:

\begin{center}
\resizebox{\columnwidth}{!}{%
\begin{tabular}{@{}lll@{}}
\hline
Problem & Positivität & Fehlende Brücke \\
\hline
Hodge & HR-Form $\geq 0$ & $\mathrm{AP} \Rightarrow$ algebraisch \\
BSD & $\hat{h} \geq 0$ & Höheres Gross-Zagier \\
P vs NP & $K^t(w|x)$ hoch & Uniformitätsbrücke \\
\hline
\end{tabular}
}
\end{center}

\subsection{Einschränkungen und intellektuelle Ehrlichkeit}

\textbf{Diese Arbeit beweist nicht $\PP \neq \NP$.} Wir stellen dies eindeutig fest. Die Beiträge sind: (1) eine saubere Umformulierung, die identifiziert, \textit{was} gezeigt werden muss (Zeugen-Inkompressibilität), (2) eine Barrierenanalyse, die argumentiert, \textit{warum} dieser Ansatz nicht von bekannten Meta-Theoremen blockiert wird, und (3) konkrete Forschungsrichtungen (Brückenstrategien B1--B4) zum Schließen der verbleibenden Lücke.

\begin{enumerate}
\item \textbf{Nichtberechenbarkeit von $K$.} Kolmogorov-Komplexität ist nicht berechenbar, was bedeutet, dass die Entropieobstruktion nicht direkt als Algorithmus ,,ausgeführt'' werden kann. Dies ist ein Feature (es umgeht natürliche Beweise), aber auch eine Einschränkung (es macht das Argument nichtkonstruktiv).

\item \textbf{Die Eindeutigkeits-Zeugen-Anforderung.} Das Entropielücken-Argument ist am stärksten für Instanzen mit eindeutigen Zeugen. Für Instanzen mit vielen Zeugen könnte ein Polynomialzeit-Algorithmus einen niedrigkomplexen Zeugen wählen, selbst wenn hochkomplexe Zeugen existieren.

\item \textbf{Äquivalenz, nicht Implikation.} Die Entropische Separationsvermutung ist logisch äquivalent zu $\PP \neq \NP$ (Remark~\ref{rem:esc-equivalence}). Daher ist das No-Go-Theorem dieser Arbeit eine \textit{Umformulierung}, keine unabhängige Beweisstrategie. Der Wert liegt in der Identifikation der informationstheoretischen Struktur des Problems.

\item \textbf{Barrierenanalyse: rigoros für $K$, heuristisch für $K^t$.} Für \emph{unbeschränktes} $K$ ist die Barrierenumgehung gut motiviert: $K$ ist nicht berechenbar (besiegt natürliche Beweise), orakelunabhängig (besiegt Relativierung) und nichtalgebraisch (umgeht plausibel Algebrisierung). Die Kernresultate verwenden jedoch \emph{ressourcenbeschränktes} $K^t$, das berechenbar ist und dessen Invarianz nur bis auf polynomiale Verlangsamung gilt. Für $K^t$ wird die Natural-Proofs-Barriere nicht automatisch umgangen, und formale Nicht-Algebrisierung bleibt offen (Remark~\ref{rem:barrier-K-vs-Kt}). Die Etablierung voller Barrierenimmunität für $K^t$-basierte Argumente ist eine wichtige Forschungsrichtung.
\end{enumerate}

\subsection{Die Uniformitätsbrücke: Lemma-Ziele}
\label{sec:uniformity-bridge}

Das zentrale offene Problem dieser Arbeit ist die \emph{Uniformitätsbrücke}: die Übertragung der Entropieseparation von unbeschränkter Kolmogorov-Komplexität $K$ (wo Barrierenimmunität gilt) auf ressourcenbeschränkte Komplexität $K^t$ (wo die berechnungstheoretische Relevanz liegt). Wir identifizieren drei konkrete Lemma-Ziele. In einer hinreichend starken und uniformen Form würde jedes einzelne einer bedingten Brückenschließung gleichkommen; in der hier formulierten Form sind es Ziele und Diagnosebegriffe, keine bewiesenen Reduktionen.

\paragraph*{Brückenziel 1: Instanzkompression.}
\emph{Aussage:} Wenn kein Polynomialzeit-Algorithmus erfüllende Belegungen von $\varphi_n$ auf weniger als $|\varphi_n|^{1-\varepsilon}$ Bits komprimieren kann, dann gilt $K^{n^c}(w | \varphi_n) \geq |\varphi_n|^{1-\varepsilon}$, oder in der stärksten Eindeutig-Zeugen-Form $K^{n^c}(w | \varphi_n) \geq |w| - O(\log n)$, für jeden gültigen Zeugen $w$.

\emph{Warum dies genügt:} Instanzkompressionshärte ist ein gut untersuchter Begriff in der parametrisierten Komplexitätstheorie~\cite{FortnowSanthanam2011,DellMelkebeek2014}. Der fehlende Schritt ist ein echtes Zeugen-zu-Kompression-Transferlemma: Ein kurzes Programm, das $w$ aus $\varphi_n$ berechnet, müsste in eine kurze, vom Verifizierer unabhängige Darstellung im präzisen Sinn von Kernelisierung oder OR-Kompressionsuntergrenzen übersetzt werden. Ein solcher Transfer folgt nicht automatisch aus der bloßen Existenz eines Zeugen-Selektors; er ist der Inhalt dieses Brückenziels.

\paragraph*{Brückenziel 2: Beweiskomplexität.}
\emph{Aussage:} Für jedes propositionslogische Beweissystem $P$ und jede Tautologie $\tau_n$, die die Unerfüllbarkeit von $\bar{\varphi}_n$ kodiert, erfüllt der Beweis $\pi$ die Bedingung $K(\pi | \tau_n) \geq |\pi|^{1-\varepsilon}$.

\emph{Warum dies genügt:} Hohe Beweiskomplexität (superpolynomiale Beweislänge) impliziert, dass Beweise selbst hohe Entropie tragen. Kombiniert mit der Entropischen Separationsvermutung ergibt sich: keine Polynomialzeit-Beweissuche kann uniform Beweise finden, weil die Beweisobjekte inkompressibel sind. Dieser Weg verbindet sich mit der Razborov--Rudich-Barriere natürlicher Beweise auf konstruktive Weise.

\paragraph*{Brückenziel 3: PRG-basierte Familien.}
\emph{Aussage:} Es existiert eine explizite Familie $\{\varphi_n\}$ von SAT-Instanzen, sodass (i) jedes $\varphi_n$ eine eindeutige erfüllende Belegung $w_n$ besitzt, (ii) $w_n$ die Ausgabe eines Pseudozufallsgenerators $G: \{0,1\}^{n/2} \to \{0,1\}^n$ ist, und (iii) $K^{n^c}(w_n | \varphi_n) \geq n/2 - O(\log n)$ gilt.

\emph{Warum dies genügt:} Wenn $G$ ein sicherer PRG ist (z.B.\ basierend auf Faktorisierungs- oder Gitterannahmen), dann folgt (iii) aus der Pseudozufälligkeit: jede effiziente Kompression von $w_n$ würde $G$ brechen. Die Eindeutig-Zeugen-Eigenschaft (i) stellt sicher, dass die Entropielücke greift (Corollary~\ref{cor:gap}). Dieser Weg reduziert die Uniformitätsbrücke auf eine kryptographische Annahme.

\begin{remark}[Minimale hinreichende Brücke]
\label{rem:minimal-bridge}
Jedes \emph{einzelne} der drei Ziele genügt nur dann, wenn es in der oben starken Form bewiesen wird. Ziel~3 reduziert auf kryptographische Standardannahmen, die selbst unbewiesen sind. Ziel~1 braucht ein neues Zeugen-zu-Kompression-Transferlemma, bevor vorhandene Instanzkompressionsuntergrenzen~\cite{Drucker2015} anwendbar werden. Ziel~2 ist das ambitionierteste und würde zugleich weit offene Fragen der Beweiskomplexität lösen.
\end{remark}

\begin{remark}[Instanzkompression als primäre Brückenstrategie]
\label{rem:compression-bridge}
Wir vertiefen Brückenziel~1 als vielversprechendste Route, dem ,,No sparsification unless PH collapses''-Programm von Fortnow--Santhanam~\cite{FortnowSanthanam2011} und Dell--van~Melkebeek~\cite{DellMelkebeek2014} folgend.

Das Argument verläuft in drei Schritten:

\textbf{Schritt A: Zeugenfinder $\Rightarrow$ niedrige bedingte Beschreibung.}
Wenn ein Polynomialzeit-Algorithmus $A$ einen gültigen Zeugen $w=A(x)$ für $x\in\mathrm{SAT}$ erzeugt, dann hat der ausgewählte Zeuge relativ zu $x$ eine niedrige bedingte Beschreibung: $K^{q_A}(w\mid x)\leq |A|+O(\log |x|)$. Das ist die bereits bewiesene einfache Richtung der Zeugen-Entropie-Charakterisierung. Es ist \emph{noch keine} Instanzkompression von $x$ im Sinn von Fortnow--Santhanam oder Dell--van~Melkebeek. Dafür bräuchte man zusätzlich eine Abbildung, die Solver-Spur, Zeugeninformation oder Verifiziererzertifikat in eine kürzere Instanz oder einen Kernel übersetzt, dessen Dekompressor außer durch das kodierte Objekt nicht von der ursprünglichen Instanz abhängt.

\textbf{Schritt B: Kompressions-No-Go $\Rightarrow$ Uniformität.}
Die bekannten Instanzkompressionsunmöglichkeitsresultate~\cite{FortnowSanthanam2011,DellMelkebeek2014,Drucker2015} zeigen: wenn SAT-Instanzen der Länge $n$ im passenden Kernelisierungs- oder OR-Kompressionssinn auf $n^{1-\varepsilon}$ Bits komprimiert werden können, kollabiert die Polynomhierarchie (NP $\subseteq$ coNP/poly). Diese Resultate dürfen hier erst eingesetzt werden, wenn Schritt~A zu einem echten Kompressionstransferlemma verstärkt wurde. Ohne diese Verstärkung ist ein Zeugen-Selektor nur ein Solver; die Kompressionsuntergrenze würde sonst Entscheidungs-Kompression mit Zeugenproduktion vermengen.

\textbf{Schritt C: Übersetzung nach $K^t$.}
Wenn ein solches Kompressionstransferlemma bewiesen wäre und die zugehörige Kompression unmöglich ist (unter PH-Nichtkollaps), dann müsste für unendlich viele Instanzen $x$ jede Polynomialzeit-Zeugenproduktion im logarithmischen Beschreibungsregime scheitern. Ein zusätzliches quantitatives Argument wäre weiterhin nötig, um dies zu $K^{n^c}(w | x) \geq |w| - O(\log n)$ für alle Zeugen $w$ zu verstärken. Diese quantitative Verstärkung ist die Nahe-Maximal-Brücke, keine Folgerung aus der Basisumformulierung.

Die verbleibenden Lücken sind also zweifach: zuerst der Zeugen-zu-Kompression-Transfer selbst, danach die Quantorenreihenfolge. Selbst ein erfolgreicher Schritt~B gäbe typischerweise ,,für jeden $A$ existieren schwere Instanzen'' (nicht-uniform), während die Uniformitätsbrücke ,,es existieren Instanzen, die für alle $A$ schwer sind'' (uniform) benötigt. Das Schließen dieser Lücke erfordert entweder (i) eine \emph{universelle} Kompressionsunmöglichkeit (nicht nur für jeden Kompressor einzeln), oder (ii) ein Diagonalisierungsargument, das über alle Polynomialzeit-Algorithmen gleichzeitig uniformisiert.
\end{remark}

\begin{remark}[Die Uniformitätsbrücke als Mischungs-Obstruktion]
\label{rem:mixing-obstruction}
Die Uniformitätsbrücke lässt eine Umformulierung in der Sprache der statistischen Mechanik zu, die direkt an das Yang--Mills-Begleitpaper anknüpft. Man betrachte das Gibbs-Maß über Zeugen, bedingt auf eine Instanz $x$:
\[
\pi_x(w) \;\propto\; \mathbf{1}\{V(x,w) = 1\} \cdot e^{-\beta H_x(w)},
\]
wobei $H_x$ eine beliebige Energiefunktion auf dem Zeugenraum ist (z.B.\ $H_x \equiv 0$ für das Gleichverteilungsmaß über gültige Zeugen). Eingeschränkte lokale, Random-Walk-, Flip- oder DPLL-Zustandsalgorithmen lassen sich manchmal als \emph{lokale Markov-Dynamik} (Glauber-/Heat-Bath-Typ) auf einem Zeugen-Schattenraum modellieren. Ein allgemeiner Polynomialzeit-Algorithmus muss keine solche Darstellung besitzen; ihn als solche sichtbar zu machen, wäre selbst eine Form der Uniformitätsbrücke.

Für diese eingeschränkten Dynamiken kann man folgendes Diagnoseziel formulieren:
\begin{quote}
\emph{Für jede NP-vollständige Sprache $L$ und jede lokale Markov-Dynamik auf Zeugen existieren Instanzfamilien $\{x_n\}$, für die die Mischungszeit zum Zeugenträger $\geq \exp(\Omega(n))$ beträgt.}
\end{quote}
Dies ist eine \textbf{Dobrushin-artige Obstruktion}: Die Einflussmatrix $C = (c_{ij})$ des Gibbs-Maßes auf dem Zeugenraum hat in der Modellanalogie Spektralradius $\rho(C) \geq 1$ für schwere Instanzen, was Dobrushin-Eindeutigkeit und damit schnelles Mischen verhindern würde. Die Parallele zum Yang--Mills-Paper~\cite{Geiger2026-YM} ist nur formal: Dort \emph{gilt} $\rho(C) < 1$ für $\beta < \beta_0$ (starke Kopplung), was eine volumenunabhängige Maßenlücke liefert; hier wird $\rho(C) \geq 1$ für NP-schwere Instanzfamilien lediglich als Suchraumdiagnose vermutet.
\end{remark}

\begin{remark}[Algorithmische Downhill-Bedingung]
\label{rem:algorithmic-downhill}
In Analogie zur Downhill-Flussbedingung (DFC) im Turbulenz-Begleitpaper~\cite{Geiger2026-TU} definiere man die \emph{Algorithmische Downhill-Bedingung} (ADC): Ein Suchprozess erzeugt eine Trajektorie $w^{(0)}, w^{(1)}, \ldots, w^{(T)}$ partieller Zeugen, und ein KL-basiertes Funktional
\[
F_x(w^{(k)}) \;=\; D_{\mathrm{KL}}\!\left(\delta_{w^{(k)}} \,\|\, \pi_x\right)
\]
misst die ,,Freie-Energie-Distanz'' zur Zeugenverteilung. Die ADC besagt, dass $F_x$ entlang jeder Polynomialzeit-Trajektorie monoton nicht-ansteigend ist. Für NP-vollständige Instanzen muss die ADC \emph{verletzt} werden: Der Suchprozess erfordert ,,Aufwärtsschritte'' (lokale Entropiezunahme), analog zum Rückstreuung in turbulenten Kaskaden. Die Anzahl der erforderlichen Aufwärtsschritte liefert eine untere Schranke für die Berechnungskomplexität der Suche -- und verbindet entropische Separation mit dem DFC-Framework.
\end{remark}

\begin{remark}[Verbindung zur Overlap Gap Property]
\label{rem:ogp-connection}
Die Mischungs-Obstruktion aus Bemerkung~\ref{rem:mixing-obstruction} hat ein präzises Gegenstück in der Literatur zur kombinatorischen Optimierung: die \emph{Overlap Gap Property} (OGP), eingeführt von Gamarnik und Sudan~\cite{GamarnikSudan2017} und zu einer systematischen Theorie algorithmischer Barrieren entwickelt~\cite{Gamarnik2021}.

Die OGP besagt, dass für zufällige Constraint-Satisfaction-Probleme nahe der Erfüllbarkeitsschwelle der Lösungsraum \emph{topologische Diskonnektivität} aufweist: Der Overlap (normierte Hamming-Distanz) zwischen zwei beliebigen nahezu optimalen Lösungen ist entweder sehr klein oder sehr groß, ohne Zwischenwerte. Dies ist ein geometrisches Analogon der Dobrushin-Obstruktion $\rho(C) \geq 1$ aus Bemerkung~\ref{rem:mixing-obstruction}: Das Gibbs-Maß über Zeugen kann in exponentiell viele wohlseparierte Cluster fragmentieren (der \emph{Shattering-Übergang} von Achlioptas und Coja-Oghlan~\cite{AchlioptasCO2008}) und breite Klassen stabiler oder lokaler Dynamiken blockieren.

Die Implikationskette lautet:
\[
\begin{aligned}
\text{Shattering}
&\Rightarrow \text{OGP}\\
&\Rightarrow \text{Versagen stabiler Alg.}\\
&\stackrel{?}{\Rightarrow} \text{hohes } K^t(w \mid x).
\end{aligned}
\]
Die ersten beiden Pfeile sind nur für eingeschränkte Algorithmenmodelle etabliert: Gamarnik, Jagannath und Wein~\cite{GamarnikJW2020} beweisen OGP-Barrieren für niedriggradige Polynome, niedrigstufige/stabile boolesche Modelle und Langevin-artige Dynamiken; Bresler und Huang~\cite{BreslerHuang2022} etablieren mithilfe der Overlap Gap Property scharfe Phasenübergänge für lokale und Message-Passing-artige Algorithmen auf zufälligem $k$-SAT.

Der dritte Pfeil -- \textbf{OGP $\Rightarrow$ hohe bedingte Kolmogorov-Komplexität $K^t(w|x)$} -- ist daher kein Theorem und in dieser direkten Allgemeinheit vermutlich falsch. Tragfähig ist ein engeres Ziel: Zeige, dass jeder Algorithmus einer vorregistrierten sichtbaren Klasse (stabil, lokal, niedriggradig, DPLL-Zustand oder Front-Filtration) einen OGP-Flaschenhals überqueren muss, und miss danach die residuale Beschreibungslänge nach Abzug von Klassentemplate und erlaubtem Rat. Der Übergang von solchen eingeschränkten sichtbaren Klassen zu beliebigen uniformen Polynomialzeit-Zeugenproduzenten ist das separate \emph{Front-Shadow}- oder \emph{Flow-Shadow}-Lemma; dieses Lemma ist die eigentliche Uniformitätsbrücke.
\end{remark}

\subsubsection*{Schwellen-Axiom und diagnostisches Framework}

Die Uniformitätsbrücke lässt eine präzise Formulierung als \emph{Schwellen-Axiom} zu -- die minimale strukturelle Bedingung, die existenzielle Härte (,,für jedes $n$ existieren schwere Instanzen'') von adversarialer Härte (,,eine einzige Prozedur erzeugt schwere Instanzen für alle $n$'') trennt.

\begin{axiom}[Uniformitätsbrücke -- UB-PvNP]\label{ax:ub-pvnp}
Sei $\mathcal{W}_n$ die Menge der Kandidaten für Hoch-Entropie-Zeugen bei SAT-Instanzen der Größe $n$: Paare $(\varphi_n, w)$ mit $K^t(w \mid \varphi_n) \geq n^\varepsilon$. Ein nahe-maximales, nicht-uniformes Brückenziel würde besagen:
\[
\begin{aligned}
&\forall\, \text{Algorithmus } A,\quad
  \exists^\infty\, n,\quad
  \exists\, \varphi_n \in \mathrm{SAT}_n:\\
&\qquad \text{alle Zeugen } w \text{ von } \varphi_n\\
&\qquad \text{erfüllen } K^t(w \mid \varphi_n) \geq n^\varepsilon.
\end{aligned}
\]
Theorem~\ref{thm:high_entropy} beweist nur die schwächere logarithmische Schwellenversion für jedes feste $(q,C)$; die angezeigte $n^\varepsilon$-Schwelle ist ein Brückenziel.
Die \textbf{Uniformitätsbrücke} verlangt eine \emph{uniforme} Konstruktion: eine einzige deterministische Prozedur $\mathcal{U}$, die bei Eingabe $1^n$ eine Instanz $\varphi_n$ mit obiger Eigenschaft erzeugt, unter Verwendung polynomieller Ressourcen in $n$.
\end{axiom}

\begin{remark}[Kernintuition]\label{rem:ub-intuition}
Nicht-Uniformität ist ,,versteckter Rat'': die schweren Instanzen $\varphi_n$ existieren, aber ihre Identität hängt in nicht-konstruktiver Weise vom Algorithmus $A$ ab. Die Uniformitätsbrücke ist der Versuch, \emph{diesen Rat zu eliminieren}. Die Härte von P\,vs\,NP liegt präzise an dieser Schwelle: wenn UB gelingt, folgt P$\,\neq\,$NP; wenn UB scheitert, muss sich das Scheitern als explizite \emph{Rat-Anforderung} manifestieren -- eine quantifizierbare Nicht-Uniformitäts-Barriere.
\end{remark}

\begin{remark}[Kandidaten-Korrespondenzen um UB-PvNP]\label{prop:ub-equivalences}
Die folgenden Punkte sind Kandidatenformulierungen und benachbarte Ziele für die Uniformitätsbrücke:
\begin{enumerate}
\item[(i)] \textbf{Uniforme Zeugenkonstruktion:} Es existiert ein Polynomialzeit-Verfahren $\mathcal{U}$, das Hoch-Entropie-Instanzen $\varphi_n = \mathcal{U}(1^n)$ für alle $n$ erzeugt.
\item[(ii)] \textbf{Uniforme Wahrheitstabellen-/Schaltkreisuntergrenzen:} Es existiert eine explizite Boolesche Funktionenfamilie $(f_n)_{n \geq 1}$, konstruierbar in $\mathrm{P}$, mit $K^t(f_n) \geq n^\varepsilon$; in der direkten Beschreibungsbuchhaltung schließt dies Schaltkreisbeschreibungen mit weniger als $n^\varepsilon$ Bits aus, also Schaltkreise der Größe $o(n^\varepsilon/\log n)$ bis auf Kodierungskonstanten.
\item[(iii)] \textbf{Uniforme Pseudozufallsgeneratoren:} Es existiert ein PRG $G: \{0,1\}^{n^\delta} \to \{0,1\}^n$, berechenbar in $\mathrm{P}$, der alle polynomialgroßen Schaltkreise täuscht.
\item[(iv)] \textbf{Parameter-uniforme Separation:} Die Konstanten in der Entropielücke $K^t(w \mid \varphi_n) \geq n^\varepsilon$ können \emph{unabhängig} von $n$ gewählt werden: keine $n$-abhängigen Schwellenwerte, Fallunterscheidungen oder Nachschlagetabellen.
\end{enumerate}
Jede der Formulierungen (i)--(iii) würde, falls sie in hinreichend starker und uniformer Form bewiesen wird, P$\,\neq\,$NP implizieren. Die Pfeile zwischen ihnen werden hier nicht als bewiesene Äquivalenzen verwendet; sie markieren Standardrouten über MCSP und Härte-gegen-Zufall, jeweils mit eigenen Hypothesen und quantitativen Verlusten.
\end{remark}

\begin{remark}[Der $n$-Leak-Scanner: Diagnostik für Nicht-Uniformität]\label{rem:n-leak}
Um zu lokalisieren, wo Uniformität bricht, durchsuche den Beweis nach diesen vier Signaturen:
\begin{enumerate}
\item \textbf{Rat-Stellen:} Jeder Schritt der Form ,,wähle für jedes $n$\ldots'' ohne eine effektive Auswahlregel.
\item \textbf{Nachschlagetabellen:} Implizite Verwendung einer $n$-spezifischen Liste oder Enzyklopädie (selbst wenn sie ,,lediglich existiert'').
\item \textbf{Diagonalisierung ohne Uniformisierung:} Ein Objekt, das durch ein reines Existenzargument (Cantor/König) definiert ist und keine uniforme Konstruktion liefert.
\item \textbf{Kodierte Instanzfamilien:} Härte, die in eine speziell vorbereitete Familie verschoben wird, die bereits $n$-abhängigen Rat enthält.
\end{enumerate}
Im aktuellen ESC-Framework sitzt der $n$-Leak in Theorem~\ref{thm:slice_nogo}: die schwere Instanz $\varphi_n$ wird durch Widerspruch erzeugt (,,wenn alle Instanzen einfach wären, würde $A$ SAT lösen''), was Existenz ohne Konstruktion liefert. Brückenziel~B3 (PRG-basierte Familien) ist der direkteste Versuch, diesen Leak zu schließen.
\end{remark}

\begin{remark}[No-Go-Mechanismen als Härtelokalisierung]\label{rem:ub-nogo}
UB-PvNP scheitert genau dann, wenn eine der folgenden Obstruktionen vorliegt:
\begin{enumerate}
\item \textbf{Rat-Unvermeidbarkeit:} Jede Kandidaten-Uniformisierung benötigt implizit $a(n)$ Bits Rat, und $a(n) \geq n^\delta$ für ein $\delta > 0$. Dies ist die P/poly-Barriere: uniforme Härte erfordert das Entkommen aus nicht-uniformer Berechnung.
\item \textbf{Parameter-Explosion:} Jede uniformisierte Konstruktion erzwingt, dass Parameter superpolynomiell in $n$ wachsen, was die ,,uniforme'' Prozedur effektiv nicht-uniform macht.
\item \textbf{Instanzfamilien-Trick:} Der Beweis funktioniert nur auf einer $n$-abhängig gewählten Familie, sodass die ,,Lösung'' in Wirklichkeit ein Nicht-Uniformitäts-Outsourcing ist.
\end{enumerate}
Jeder No-Go-Mechanismus würde, wenn bewiesen, ein \emph{bedingtes Unmöglichkeitstheorem} darstellen: ,,P$\,\neq\,$NP kann nicht mit Methode $X$ bewiesen werden, weil $X$ $a(n)$ Bits Rat benötigt.'' Solche Resultate sind Barrierentheoreme im Geiste von Baker--Gill--Solovay und Razborov--Rudich, aber gezielt auf das spezifische Uniformitätsdefizit.
\end{remark}

\begin{remark}[Advice-Erosion-Kurve]\label{rem:advice-erosion}
Der $n$-Leak aus Bemerkung~\ref{rem:n-leak} lässt sich als explizites
Ratsbudget quantifizieren.  Für eine samplbare Ja-Verteilung
$\mathcal{D}_n$ definiere man als informelles Diagnoseprofil
\[
E_{\mathcal{D}}(n,b,T)
  := \inf_{A:\,|a_n|\leq b}
     H_{\mathrm{res}}\bigl(W \mid X,A(X,a_n),T\bigr),
\]
wobei $A$ über $T$-Zeit-Verfahren mit höchstens $b$ Bits
$n$-abhängigem Rat läuft.  Die Lesart ist: $b=0$ bedeutet echte
Uniformität, $b=O(\log n)$ kodiert nur Größen- oder Regimeinformation,
während $b\geq n^\delta$ versteckte nicht-uniforme Adressinformation
darstellt.

Die sichere einfache Richtung ist nur diagnostisch: Hat ein NP-Suchproblem
einen uniformen polynomialzeitlichen Zeugen-Selektor, dann kollabiert das
zugehörige Advice-Erosion-Profil bereits bei $b=0$ und $T=\poly(n)$ für
die vom Selektor induzierte Zeugenwahl.  Das offene Brückenziel ist der
Gegen-Stresstest: SAT-/MCSP-/OGP-Familien zu finden, für die
$E_{\mathcal{D}}(n,\operatorname{polylog}(n),\poly(n))$ polynomial groß bleibt.  Das würde
für sich allein P$\,\neq\,$NP nicht beweisen, aber exakt lokalisieren, wie
viel nicht-uniforme Information nötig ist, bevor die Entropieobstruktion
erodiert.
\end{remark}

\begin{remark}[Search-Flow-Conductance]\label{rem:search-flow-conductance}
Eine komplementäre Diagnose misst Transport statt Rat.  Für eine
SAT-Instanz $x$ sei $\Omega_x$ ein kanonischer Zustandsraum aus
Zeugen-Schatten und $S_x\subseteq\Omega_x$ die Region gültiger Zeugen.
Eine eingeschränkte Suchklasse $\mathcal{C}$ (etwa lokale, Markovsche,
DPLL-Zustands- oder Flip-Prozess-Modelle) induziert Verteilungen
\[
\mu_0 \to \mu_1 \to \cdots \to \mu_T
\]
auf $\Omega_x$.  Informell ist $\Phi_{\mathcal{C}}(x)$ die kleinste
Flusskapazität über einen Schnitt, der die Startmasse von $S_x$ trennt.

Für solche eingeschränkten Modelle kann man ein normales
Flaschenhals-Lemma anstreben: Ist $\mu_0(S_x)$ vernachlässigbar und hat
jeder trennende Schnitt Kapazität
$\Phi_{\mathcal{C}}(x)\leq \exp(-n^\varepsilon)$, dann kann kein
$\poly(n)$-Schritt-Prozess in $\mathcal{C}$ konstante Maße auf $S_x$
legen.  Das ist nur eine Obstruktion für eingeschränkte Modelle.  Das
wirklich harte ``Flow-Shadow''-Lemma müsste zeigen, dass jeder uniforme
polynomialzeitliche Zeugenproduzent für die gewählte SAT-Familie ohne
superpolynomialen Verlust eine solche kanonische Flussdarstellung besitzt.
Dieses Flow-Shadow-Lemma ist eine weitere präzise Form der
Uniformitätsbrücke.
\end{remark}

\begin{remark}[Capillary-Front-Filtration]\label{rem:capillary-front}
Die Advice- und Flow-Diagnosen lassen sich zu einem strengeren
Front-Ledger verbinden.  Für einen Algorithmus $A$, eine Instanz $x$ und
Rat $a_n$ betrachte man eine vorregistrierte Folge
\[
F_0(A,x,a_n)\leq F_1(A,x,a_n)\leq\cdots\leq F_T(A,x,a_n)
\]
niedrig komplexer sichtbarer Fronten in einem Zeugen-Schattenraum:
partielle Belegungen, DPLL-Zustandsklassen, Clusterzertifikate,
Resolution-Breitenbänder oder Bounded-Proof-Family-Templates.  Ein
Kapillartreffer ist ein isolierter Pfad zu einem Zeugen; eine Front ist
eine niedrig komplexe Struktur, die positive Zeugenmasse erreicht.

Zwei Diagnosegrößen müssen getrennt bleiben.  Der Shannon-Proxy ist
\[
H_{\mathrm{front}}(x,t)=H_{\mathrm{res}}(W\mid x,F_t),
\]
während das eigentliche ressourcenbeschränkte/MDL-Residuum für einen Zeugen
$w$ lautet:
\[
\begin{aligned}
K_{\mathrm{front}}^t(x,w,t)
  &=K^t(w\mid x)-K(\mathrm{BPF})\\
  &\quad -K(\mathrm{Template})-|a_n|-K(r\mid x),
\end{aligned}
\]
wobei $r$ einen erlaubten Spezialisierungsparameter der sichtbaren Front
bezeichnet.  Die erste Größe ist nur experimentelle Evidenz; eine
Uniformitätsobstruktion verlangt, dass die zweite Größe nach Abzug von
Template-, Rat- und Spezialisierungsbits groß bleibt.

Für eine eingeschränkte Klasse $\mathcal{C}$ könnte ein Lemma lauten: Wenn
jede erlaubte Front in $\poly(n)$ Zeit entweder nur $o(1)$ Masse von $S_x$
trifft oder auf den residualen Zeugen
$K_{\mathrm{front}}^t(x,w,t)\geq n^\varepsilon$ übrig lässt, dann kann
$\mathcal{C}$ die Familie nicht mit konstanter Erfolgswahrscheinlichkeit
lösen.  Dies bleibt eine Aussage über eingeschränkte Modelle.  Um sie gegen
alle Polynomialzeit-Algorithmen zu verwenden, müsste ein Front-Shadow-Lemma
zeigen, dass jeder uniforme Zeugenproduzent ohne superpolynomialen Verlust
eine solche niedrig komplexe sichtbare Front induziert.

Das Leakage-Ledger muss außerdem ausweisen, woher die Front stammt.  Eine
Front, die aus demselben Solverlauf extrahiert wird, dessen Zeugenmasse sie
erklären soll, ist kein externes Uniformitätssignal, sofern die Extraktionsregel
nicht vollständig als Template- oder Ratlänge verbucht wird.  Sichere
Buchhaltung liefert die elementare Schranke
\[
K^{q_T}(w\mid x)\leq K(\mathcal{C}_n)+|a_n|+K(r\mid x)+O(\log n),
\]
für jede vorregistrierte sichtbare Produzentenklasse $\mathcal{C}_n$ mit Rat
$a_n$, Spezialisierungsparameter $r$ und Laufzeit $T(n)$.  Ein positives
$K^t$-Residuum schließt daher nur die verbuchte sichtbare Klasse aus, bis ein
separates Front-Shadow- oder Flow-Shadow-Lemma die Buchhaltung auf beliebige
uniforme polynomialzeitliche Zeugenproduzenten überträgt.
\end{remark}

\begin{lemma}[UB impliziert Trennung]\label{lem:ub-separation}
Wenn UB-PvNP (Axiom~\ref{ax:ub-pvnp}) mit einer Zeitschranke gilt, die jeden
festen polynomialzeitlichen Zeugen-Selektor dominiert, dann:
\begin{enumerate}
\item[(a)] P $\neq$ NP (unmittelbar aus der Existenz uniformer schwerer Instanzen).
\item[(b)] die von $\mathcal{U}$ erzeugte explizite Familie besitzt die
behauptete uniforme Zeugenentropie-Obstruktion:
\[
  K^t(w\mid \varphi_n)\geq n^\varepsilon
\]
für alle Zeugen von $\varphi_n$ auf unendlich vielen $n$.
\end{enumerate}
UB-PvNP allein beweist weder die Existenz von Einwegfunktionen noch eine
globale Slice-Entropy-Untergrenze für SAT-Wahrheitstabellen.  Das sind
separate Brückenziele: Einwegfunktionen benötigen eine zusätzliche
PRG-/Average-Case-$K^t$-Route (etwa über Liu--Pass oder passende
Hardness-vs-Randomness-Hypothesen), während Slice-Entropy-Untergrenzen einen
Wahrheitstabellen- oder Schaltkreis-Schranken-Transfer benötigen.
Umgekehrt impliziert P $\neq$ NP genau dann UB-PvNP, wenn die schweren Instanzen \emph{explizit} gemacht werden können -- was genau der Inhalt der Uniformitätsbrücke ist.
\end{lemma}

\subsection{Das Meta-Muster}

Über alle Probleme in diesem Forschungsprogramm hinweg ist die Struktur:

\begin{enumerate}
\item \textbf{Identifiziere die positive Form} (Höhe, HR-Form, Kapazität, Entropie).
\item \textbf{Zeige eine Richtung} (algebraisch $\Rightarrow$ positiv, $\PP \Rightarrow$ niedrige Entropie).
\item \textbf{Zeige, dass die schwere Richtung für Obstruktionen versagt} (nichtalgebraisch verletzt Positivität, NP-schwer hat hohe Entropie).
\item \textbf{Schließe die Schleife mit einem Kompositions-/Kopplungstheorem.}
\end{enumerate}

Für P vs NP sind Schritte 1--3 in dieser Arbeit etabliert. Schritt 4 -- die Uniformitätsbrücke -- bleibt die zentrale Herausforderung.

\subsection{Offene Richtungen}

\begin{remark}[Levins universelle Suche und Uniformitätskollaps]
\label{rem:levin-uniformity}
Levins universeller Suchalgorithmus $U$ führt alle Programme der Länge $\leq l$ in verschränkter Weise aus und weist Programm $p$ die Zeit $2^{-|p|} \cdot T$ zu. Wenn P $=$ NP, dann findet $U$ Zeugen in Polynomialzeit (mit einer großen Konstante). Die Nahe-Maximal-Zeugenentropie-Vermutung liefert eine komplementäre Perspektive: wenn die relevanten Zeugen $K^t(w|\varphi) \geq |w| - O(\log n)$ erfüllen, dann wird $U$'s Zeitverteilung vom \emph{korrekten} Programm dominiert, dessen Länge im Wesentlichen $|w|$ ist. Die resultierende Laufzeit ist $2^{|w|} \cdot \mathrm{poly}(n)$ -- exponentiell. Dies ist ein heuristisches Optimalitätsbild für Levin-Suche unter der stärkeren Nahe-Maximal-Hypothese, keine Folge von Theorem~\ref{thm:high_entropy}.
\end{remark}

\begin{remark}[Valiant-Vazirani-Reduktion und Wenig-Zeugen-Verstärkung]
\label{rem:valiant-vazirani}
Das Valiant-Vazirani-Theorem reduziert SAT mit hoher Wahrscheinlichkeit über zufällige Hashfunktionen auf Unique-SAT. Wenn die Nahe-Maximal-Zeugenentropie-Vermutung für die entstehenden Eindeutig-Zeugen-Instanzen gilt, dann erfüllt für einen zufälligen Hash $h$ der eindeutige Zeuge $w_h$ von $\varphi \wedge h(w) = 0^k$ mit hoher Wahrscheinlichkeit (über $h$) die Bedingung $K^t(w_h | \varphi, h) \geq |w_h| - O(\log n)$. Diese ,,Wenig-Zeugen-Verstärkung'' würde die stärkere Nahe-Maximal-Variante auf generische SAT-Instanzen übertragen, auf Kosten einer randomisierten Reduktion. Derandomisierung über Nisan-Wigderson-Generatoren wäre ein zusätzliches Brückenziel.
\end{remark}

\begin{remark}[Min-Entropie-Variante]
\label{rem:min-entropy}
Die Nahe-Maximal-$K^t$-Formulierung kann zu einer distributionellen Version mittels Min-Entropie abgeschwächt werden: für eine samplbare Verteilung $\mathcal{D}$ über SAT-Instanzen definiere $H_\infty^t(\mathcal{D}) = \min_{(\varphi, w) \in \mathrm{supp}(\mathcal{D})} K^t(w | \varphi)$. Wenn $H_\infty^t(\mathcal{D}) \geq |w| - O(\log n)$ für eine samplbare $\mathcal{D}$ gilt, dann folgt Average-Case-Härte von NP (Impagliazzos ,,Pessiland''). Diese Abschwächung ist ein distributionelles Brückenziel für die stärkere Nahe-Maximal-Variante und verbindet sich mit der reichhaltigen Theorie der Average-Case-Komplexität und Einwegfunktionen.
\end{remark}

\begin{remark}[Adversariale Zeugenfamilie]
\label{rem:adversarial-witness}
Eine stärkere Variante der ESC konstruiert eine explizite \emph{adversariale} Zeugenfamilie $\{w_n\}$ mit einer globalen Konsistenzprüfung: die Zeugen sind nicht unabhängig schwer, sondern erfüllen eine kollektive Nebenbedingung $\mathcal{C}(w_1, \ldots, w_N) = 1$, die leicht zu verifizieren, aber schwer zu erfüllen ist. Jeder Algorithmus, der $w_n$ für alle $n \leq N$ berechnet, muss die globale Struktur $\mathcal{C}$ ,,kennen'', was eine Beschreibungslänge $\Omega(N)$ erfordert. Dieser Ansatz vermeidet die Einzel-Instanz-Einschränkungen von Korollar~\ref{cor:gap} und ist mit der Theorie interaktiver Beweise und PCP verwandt.
\end{remark}

\begin{remark}[Kompressions-zu-Schaltkreis-Transfer]
\label{rem:compression-circuit}
Für eingeschränkte Schaltkreisklassen $\mathcal{C}$ (z.B.\ $\mathrm{AC}^0$, monotone Schaltkreise) nimmt der Kompressions-zu-Schaltkreis-Transfer eine schärfere Form an: wenn $\varphi_n$ keinen $\mathcal{C}$-Schaltkreis der Größe $s$ hat, der seinen Zeugen berechnet, dann gilt $K^t_{\mathcal{C}}(w | \varphi_n) \geq \log \binom{2^n}{s}$, wobei $K^t_{\mathcal{C}}$ die $\mathcal{C}$-eingeschränkte Kolmogorov-Komplexität bezeichnet. Für $\mathrm{AC}^0$ liefern die Razborov-Smolensky-Schranken $s \geq 2^{n^{\Omega(1)}}$, was superpolynomiales $K^t_{\mathrm{AC}^0}$ bedingungslos ergibt. Dies liefert einen ,,Machbarkeitsnachweis'' für Brückenziel~1 in einem eingeschränkten Modell.
\end{remark}

\begin{remark}[MCSP-Internalisierungsroute]
\label{rem:mcsp-internalisation}
Allender und Das (2017) zeigten, dass MCSP (Minimum Circuit Size Problem) unter randomisierten Reduktionen schwer für SZK ist. Liu und Pass~\cite{LiuPass2020} bewiesen, dass Average-Case-Härte von $K^t$ (ein enger Verwandter von MCSP) \emph{äquivalent} zur Existenz von Einwegfunktionen ist. Kombiniert mit Hiraharas~\cite{Hirahara2018} nicht-Black-Box Worst-Case-zu-Average-Case-Reduktionen innerhalb von NP liefert das eine Route von MCSP-Härte zu OWF-Existenz. Unser Framework erweitert diese Kette:
\[
\begin{aligned}
\text{MCSP schwer}
&\xRightarrow{\text{Liu--Pass}} \text{OWF existieren}\\
&\xRightarrow{\text{Imp.--Levin}} \text{avg-case NP schwer}\\
&\Rightarrow \PP \neq \NP .
\end{aligned}
\]
Die erste Implikation nutzt die Liu--Pass-Äquivalenz zwischen $K^t$-Average-Case-Härte und OWF-Existenz (die Verbindung zu MCSP kommt aus der Allender-et-al.-Korrespondenz zwischen MCSP und $K^t$); die zweite ist eine Standardreduktion; die dritte ist unmittelbar, denn Average-Case-NP-Härte impliziert $\PP \neq \NP$, da $\PP=\NP$ NP auch im Durchschnitt leicht machen würde. Der Beweis von Average-Case-Härte für MCSP würde daher über Liu--Pass plus Standardreduktionen $\PP \neq \NP$ liefern.
\end{remark}

\begin{remark}[Ilangos MCSP-NP-Härte]\label{rem:ilango-mcsp}
Hirahara und Ilango \cite{Ilango2025} liefern Evidenz für MCSP-Härte, indem sie NP-Härte unter deterministischen quasi-polynomiellen nicht-adaptiven Reduktionen aus gut untersuchten kryptographischen und schaltkreisbezogenen Annahmen beweisen. Dies löst die unbedingte NP-Härte von MCSP nicht. Für das vorliegende Framework ist die sicherere Relevanz enger: Das Resultat identifiziert MCSP als plausibles Brückenobjekt, um zeitbeschränkte Kolmogorov-Komplexität in standardmäßige komplexitätstheoretische untere Schranken zu übersetzen.
\end{remark}

\begin{remark}[Quantenerweiterung]
\label{rem:quantum}
Das Entropie-Framework lässt sich natürlich auf Quantenberechnung erweitern. Ein Quantenalgorithmus mit $T$ Gattern auf $n$ Qubits kann durch $O(T \log T)$ klassische Bits beschrieben werden (Kodierung der Gattersequenz), also $K(\mathrm{desc}(U)) \leq O(T \log T)$ für die klassische Beschreibung des Unitäroperators $U$. Falls $\mathrm{BQP} = \mathrm{QMA}$, dann würde der Polynomialzeit-Quantenalgorithmus den QMA-Zeugen ersetzen, und das analoge Zeugen-Entropielücken-Argument würde gelten: die ressourcenbeschränkte Quanten-Kolmogorov-Komplexität des Zeugen müsste gleichzeitig niedrig sein (durch den BQP-Algorithmus) und hoch (für Instanzen, die zufällige Zeichenketten kodieren). Dies legt $\mathrm{BQP} \neq \mathrm{QMA}$ nahe, obwohl die Rigorgisierung ein Quanten-Analogon der ressourcenbeschränkten Kolmogorov-Komplexität erfordert, was ein aktives Forschungsgebiet ist.
\end{remark}

\begin{remark}[Spieltheoretisches Interaktionsmodell]
\label{rem:game-theoretic}
Die P-vs-NP-Frage lässt eine spieltheoretische Formulierung zu: ein Entropie-Maximierer (,,Natur'') wählt Zeugenverteilungen, um $K^t(w|\varphi)$ zu maximieren, während ein Schaltkreis-Minimierer (,,Algorithmus'') Schaltkreise wählt, um die Berechnungszeit zu minimieren. Der Minimax-Wert dieses Nullsummenspiels beschreibt den optimalen Kompromiss zwischen Zeugenentropie und Schaltkreisgröße. Eine stärkere Nahe-Maximal-Variante würde behaupten, dass die Natur diesen Wert in einem skalensensitiven Sinn für unendlich viele $\varphi$ positiv halten kann. Diese Formulierung verbindet sich mit der Verwendung von Minimax-Prinzipien im FST-Programm in anderen Domänen (Yang--Mills, Turbulenz), ist aber kein Theorem dieser Arbeit.
\end{remark}

\begin{remark}[Quanten-Zertifikate und MIP*]
\label{rem:quantum-certificates}
Die Nahe-Maximal-Zeugenentropie-Variante kann in die Sprache quanteninteraktiver Beweise übersetzt werden: Eine Entropielücke $K^t(w|\varphi) \geq |w| - O(\log n)$ würde implizieren, dass kein polynomialgroßer Quantenschaltkreis einen Quantenzustand $|\psi_w\rangle$, der den Zeugen kodiert, allein aus $|\varphi\rangle$ präparieren kann. Im MIP*-Framework (Natarajan--Wright~\cite{NatarajanWright2019}, Ji--Natarajan--Vidick--Wright--Yuen~\cite{JiNatarajanVidickWrightYuen2021}) kann Verschränkung zwischen Beweisern als Mechanismus betrachtet werden, die Zeugenentropie über nichtkommunizierende Parteien zu ,,verteilen''. Ob Verschränkung eine solche verstärkte Entropiebarriere umgehen kann, ist eine offene Analogie, kein Claim dieses Frameworks.
\end{remark}

% Remark on Kolyvagin-type hierarchy removed: the analogy between Euler systems
% (arithmetic geometry of elliptic curves, Galois cohomology, Selmer groups)
% and SAT-witness spaces is too vague to be productive. The only shared property
% -- "hardness propagates from one instance to related instances" -- is not
% specific enough to constitute a research direction.

\begin{remark}[Selbstreferentielle Perspektive -- informell]
\label{rem:self-referential}
Eine heuristische Perspektive darauf, warum das Schließen der Uniformitätsbrücke schwer ist: Für jede feste Polynomialzeit-Maschine $A$ mit Beschreibungslänge $|A|$ und Laufzeitpolynom $q_A$ liefert Theorem~\ref{thm:high_entropy} (bedingt auf $\PP \neq \NP$) für jede hinreichend große logarithmische Konstante $C$ unendlich viele Instanzen $\varphi$, deren gültige Zeugen alle $K^{q_A}(w \mid \varphi) > C\log|\varphi|$ erfüllen. Da $A$ ein festes Programm ist, gilt jedoch $K^{q_A}(A(x) \mid x) \leq |A| + O(\log |x|) \leq C\log |x|$ für alle großen $x$. Diese beiden Schranken sind auf den entsprechenden harten Instanzen inkompatibel: $A$ kann dort keinen gültigen Zeugen produzieren. Dies ist lediglich die nicht-uniforme Richtung der Zeugen-Entropielücke (bereits in Theorem~\ref{thm:high_entropy} etabliert) und schließt die Uniformitätsbrücke \emph{nicht}; dafür braucht man harte Instanzen, die allen Algorithmen gleichzeitig widerstehen, nicht nur jedem Algorithmus einzeln.
\end{remark}

\begin{remark}[Extended Frege und MCSP]
\label{rem:frege-mcsp}
Zwei konkrete Ziele für bedingungslosen Fortschritt:
\begin{enumerate}
\item \emph{Extended Frege:} Wenn superpolynomiale untere Schranken für die Extended-Frege-Beweislänge für Tautologien etabliert werden können, die die Unerfüllbarkeit ESC-schwerer Instanzen kodieren, dann haben die Beweisobjekte selbst hohe Entropie ($K(\pi | \tau) \geq |\pi|^{1-\varepsilon}$), was Brückenziel~2 realisiert.
\item \emph{MCSP und Nichtuniformität:} Das Minimum Circuit Size Problem (Allender--Hirahara~\cite{AllenderHirahara2019}) ist die natürliche Entscheidungsversion der Kolmogorov-Komplexität für Schaltkreise. MCSP-Untergrenzen sind eine benachbarte Route, um Kompressionsaussagen in Schaltkreisuntergrenzen zu übersetzen; $\mathrm{MCSP}\notin\PP/\poly$ würde EH aber nicht für sich allein implizieren, solange keine zusätzliche Reduktion von kleinen SAT-Schnittbeschreibungen auf kleine MCSP-Schaltkreise oder ein vergleichbares uniformes Kompressions-zu-Schaltkreis-Theorem vorliegt.
\end{enumerate}
\end{remark}

\begin{remark}[Analogie zur Hodge-Vermutung]
\label{rem:hodge-analogy}
Auf einer sehr hohen Ebene teilt das ESC-Framework ein Strukturmuster mit der Hodge-Vermutung: Objekte mit strukturierten Beschreibungen (Polynomialzeit-Algorithmen / algebraische Zyklen) bilden eine Teilmenge eines größeren Raums, und die Frage ist, ob bestimmte ausgezeichnete Objekte (NP-Zeugen / Hodge-Klassen) in dieser Teilmenge liegen. Diese Analogie ist jedoch \emph{oberflächlich}: Die Hodge-Riemann-Bilinearform ist ein tiefer topologischer Invariant auf Kohomologie, während $K^t \geq 0$ eine triviale Eigenschaft jeder Längenfunktion ist. Die Analogie ist als heuristische Parallele in der Beweisstrategie zu verstehen, nicht als Evidenz für eine mathematische Verbindung.
\end{remark}


% ===================================================================
\subsection{Post-Zookeeper-Transfer: Williams-normalisierte Uniformität}
\label{subsec:zookeeper-pnp-transfer}

Die Post-Zookeeper-Rolle dieses Papers besteht darin, die
Uniformitätslücke operational zu machen.  Der relevante aktuelle Input
ist Williams' Square-Root-Space-Simulation zeitbeschränkter Berechnung:
Zeit $t$ kann im relevanten Turingmaschinenmodell mit
$O(\sqrt{t\log t})$ Speicher simuliert werden.  Das löst P versus NP
nicht, verändert aber die Ressourcengeometrie, die die
Entropieobstruktion überstehen muss.

Relativ zum aktuellen CORE-Status ist die Zookeeper-Schließung von
Riemann-$C2/MS2$ eine \emph{Normalform-Bedingung}, kein Transferbeweis
für Komplexitätstheorie.  Sie legitimiert die fünf Felder unten, aber
sie schließt die Uniformity Bridge nicht.  Der offene mathematische
Schritt dieses Papers bleibt der Übergang von nicht-uniformen
hochentropischen Zeugen zu einer uniformen unteren Schranke, die die beste
verfügbare Zeit-Speicher-Simulation übersteht.

Die Transfer-Normalform lautet:
\begin{description}
  \item[Operator.] Uniforme Such- und Verifikationsoperatoren, gemessen
  durch ressourcenbeschränkte Kolmogorov-Komplexität $K^t$.
  \item[Defekt.] Der Uniformitätsdefekt: Existenz hochentropischer
  Zeugen minus uniforme Erzeugbarkeit unter der besten verfügbaren
  Zeit-Speicher-Simulation.
  \item[Approximation.] SAT-Schnitte, Resolution-Width-Regime,
  MCSP/XOR/MUX-Kandidaten und Circuit-Size-Level.
  \item[Gap.] Ein verbleibender Entropy-Hardness-Gap nach Normalisierung
  der Zeit durch die Square-Root-Space-Simulation.
  \item[Grenzübergang.] Übergang von nicht-uniformen Instanzfamilien zu
  einer uniformen Sprache bzw. einem Algorithmus.
\end{description}

Daraus ergibt sich eine verfeinerte Diagnostik:
\[
\begin{aligned}
  H_{\mathrm{slice}}^{\mathrm{ts}}(n;s)
  &= \text{Zeugenentropie nach Simulation}\\
  &\quad \text{in Speicher }s,\qquad s \asymp \sqrt{t\log t}.
\end{aligned}
\]
Das nächste Ziel ist kein direkter Beweis von
$\mathrm{P}\ne\mathrm{NP}$, sondern ein Transferlemma: Existiert eine
uniforme polynomielle Suchprozedur, dann muss diese
Williams-normalisierte Schnittentropie entlang der zugehörigen
Ressourcenfläche kollabieren.  MCSP-artige Kandidaten und
Resolution-Width-Familien sind dann die natürlichen Stresstests gegen
Kollaps.

\subsection*{Ergebnisklassifikation}

Dieses Paper ist ein \textbf{Reformulierungspaper}: Es etabliert eine Äquivalenz zwischen der P\,$\neq$\,NP-Vermutung und einer entropischen Bedingung (ESC), ohne die Separation selbst zu beweisen.

\begin{table*}[t]
\caption{Ergebnisklassifikation.}
\label{tab:ergebnisklassifikation}
\centering
\footnotesize
\resizebox{\textwidth}{!}{%
\begin{tabular}{@{}lll@{}}
\toprule
\textbf{Ergebnis} & \textbf{Typ} & \textbf{Referenz} \\
\midrule
\multicolumn{3}{@{}l}{\textit{Bewiesene Äquivalenzen:}} \\
ESC $\Leftrightarrow$ P\,$\neq$\,NP & Proposition + Bemerkung & Prop.~\ref{prop:equiv}, Bem.~\ref{rem:esc-equivalence} \\
EH $\Rightarrow$ P\,$\neq$\,NP & Theorem & Thm.~\ref{thm:slice_nogo} \\
\midrule
\multicolumn{3}{@{}l}{\textit{Strukturelle Beiträge:}} \\
Barriere-Immunität (S1--S5 für $K$; heuristisch für $K^t$) & Analyse & Abschn.~\ref{sec:barriers}, Bem.~\ref{rem:barrier-K-vs-Kt} \\
Resolutionsbreite-Vermutung & Vermutung & Verm.~\ref{conj:resolution-width} \\
MCSP-Internalisierungsroute & Bemerkung & Bem.~\ref{rem:mcsp-internalisation} \\
Advice-Erosion-Kurve & Brückenziel & Bem.~\ref{rem:advice-erosion} \\
Search-Flow-Conductance & Brückenziel & Bem.~\ref{rem:search-flow-conductance} \\
Capillary-Front-Filtration & Brückenziel & Bem.~\ref{rem:capillary-front} \\
\midrule
\multicolumn{3}{@{}l}{\textit{Explorative Numerik (keine Evidenz im Beweissinn):}} \\
Phasenübergangsstruktur beobachtet & illustrativ & Bem.~\ref{rem:sat-entropy} \\
$H_{\mathrm{slice}}$-Wachstumstrend ($n \leq 16$, 5 Punkte) & illustrativ & Bem.~\ref{rem:sat-entropy} \\
\midrule
\multicolumn{3}{@{}l}{\textit{Offene Brückenziele:}} \\
Uniformitätsbrücke & Offenes Problem & Abschn.~\ref{sec:uniformity-bridge} \\
Entropie-Härte (EH) & Starke hinreichende Vermutung & Abschn.~\ref{sec:slices} \\
\bottomrule
\end{tabular}
}
\end{table*}

% ===================================================================
% 12. SCHLUSSFOLGERUNG
% ===================================================================
\section{Schlussfolgerung}
\label{sec:conclusion}

Die Äquivalenz zwischen $\PP \neq \NP$ und der Inkompressibilität von NP-Zeugen (unter ressourcenbeschränkter Kolmogorov-Komplexität) ist in verschiedenen Formen seit den Arbeiten von Orponen et al.~\cite{OrponenKoSchoeningWatanabe1994} und Fortnow~\cite{Fortnow2000} bekannt. Diese Arbeit präsentiert diese bekannten Verbindungen in einem einheitlichen Framework und führt die Terminologie der algorithmischen Positivität und des entropischen No-Go-Theorems ein.

Die Zeugen-Entropielücke -- die Diskrepanz zwischen der hohen bedingten Komplexität von Zeugen (unter $\PP \neq \NP$) und der niedrigen Komplexität, die Polynomialzeit-Algorithmen erzeugen können -- liefert eine saubere Charakterisierung dessen, was jeder Beweis von $\PP \neq \NP$ etablieren muss. Die Barrierenanalyse argumentiert, dass $K$-basierte Argumente nicht von den drei klassischen Barrieren blockiert werden, obwohl diese Immunität rigoros nur für unbeschränktes $K$ gilt und für das in den Kernresultaten verwendete ressourcenbeschränkte $K^t$ heuristisch bleibt.

Die verbleibende Herausforderung ist die Uniformitätsbrücke: die Konvertierung der existenziellen Entropieschranke in eine universelle untere Schranke gegen alle Polynomialzeit-Algorithmen. Die identifizierten Brückenstrategien (Instanzkompression, Beweiskomplexität, PRG-basierte Familien, MCSP) sind selbst große offene Probleme, im Kern so schwer wie P vs NP. Der Wert des Frameworks liegt darin, diese Struktur explizit zu machen.

\begin{quote}
\textit{,,P vs NP fragt nicht, ob schwere Probleme existieren. Es fragt, ob das Universum effizienter Berechnung Raum für Probleme hat, deren Lösungen irreduzible Information tragen -- und die Entropietheorie liefert die natürliche Sprache, um diese Frage mit voller Präzision zu formulieren.''}
\end{quote}


% ===================================================================
% DANKSAGUNGEN
% ===================================================================
\begin{acknowledgments}
Diese Arbeit wurde als unabhängige Forschung ohne institutionelle Förderung durchgeführt.

\textit{KI-Werkzeuge.} KI-basierte Werkzeuge (Claude, Anthropic) wurden für mathematische Formalisierung und Texterstellung verwendet. Alle Inhalte und die Verantwortung für die Korrektheit liegen ausschließlich beim Autor.

\textit{Förderinformation.} Diese Forschung erhielt keine externe Förderung.
\end{acknowledgments}


% ===================================================================
% BIBLIOGRAPHIE
% ===================================================================
\begin{thebibliography}{34}

\bibitem{BakerGillSolovay1975}
T.~Baker, J.~Gill, \& R.~Solovay (1975).
Relativizations of the $\mathrm{P} =\mathrel{?} \mathrm{NP}$ question.
\textit{SIAM Journal on Computing}, 4(4), 431--442.
DOI: 10.1137/0204037.

\bibitem{RazborovRudich1997}
A.~Razborov \& S.~Rudich (1997).
Natural proofs.
\textit{Journal of Computer and System Sciences}, 55(1), 24--35.
DOI: 10.1006/jcss.1997.1494.

\bibitem{AaronsonWigderson2009}
S.~Aaronson \& A.~Wigderson (2009).
Algebrization: A new barrier in complexity theory.
\textit{ACM Transactions on Computation Theory}, 1(1), 2.
DOI: 10.1145/1490270.1490272.

\bibitem{LiVitanyi2008}
M.~Li \& P.~Vit\'anyi (2008).
\textit{An Introduction to Kolmogorov Complexity and Its Applications}, 3rd ed.
Springer.
DOI: 10.1007/978-0-387-49820-1.

\bibitem{Cook1975}
S.~A.~Cook (1975).
Feasibly constructive proofs and the propositional calculus.
\textit{Proceedings of the 7th ACM Symposium on Theory of Computing}, pp.\ 83--97.
DOI: 10.1145/800116.803756.

\bibitem{BenSasson2001}
E.~Ben-Sasson \& A.~Wigderson (2001).
Short proofs are narrow -- resolution made simple.
\textit{Journal of the ACM}, 48(2), 149--169.
DOI: 10.1145/375827.375835.

\bibitem{Razborov1985}
A.~Razborov (1985).
Lower bounds on the monotone complexity of some Boolean functions.
\textit{Doklady Akademii Nauk SSSR}, 281(4), 798--801.

\bibitem{Razborov1987}
A.~Razborov (1987).
Lower bounds on the size of bounded depth circuits over a complete basis with logical addition.
\textit{Mathematical Notes}, 41(4), 333--338.
DOI: 10.1007/BF01137685.

\bibitem{Grigoriev2001}
D.~Grigoriev (2001).
Linear lower bound on degrees of Positivstellensatz calculus proofs for the parity.
\textit{Theoretical Computer Science}, 259(1--2), 613--622.
DOI: 10.1016/S0304-3975(00)00403-7.

\bibitem{Mulmuley2011}
K.~D.~Mulmuley (2011).
On P vs.\ NP and geometric complexity theory.
\textit{Journal of the ACM}, 58(2), 5.
DOI: 10.1145/1667053.1667055.

\bibitem{Arora2009}
S.~Arora \& B.~Barak (2009).
\textit{Computational Complexity: A Modern Approach}.
Cambridge University Press.

\bibitem{Kolmogorov1965}
A.~N.~Kolmogorov (1965).
Three approaches to the quantitative definition of information.
\textit{Problems of Information Transmission}, 1(1), 1--7.

\bibitem{Chaitin1966}
G.~J.~Chaitin (1966).
On the length of programs for computing finite binary sequences.
\textit{Journal of the ACM}, 13(4), 547--569.
DOI: 10.1145/321356.321363.

\bibitem{Monasson1999}
R.~Monasson, R.~Zecchina, S.~Kirkpatrick, B.~Selman, \& L.~Troyansky (1999).
Determining computational complexity from characteristic `phase transitions'.
\textit{Nature}, 400, 133--137.
DOI: 10.1038/22055.

\bibitem{ImpagliazzoWigderson1997}
R.~Impagliazzo \& A.~Wigderson (1997).
P = BPP if E requires exponential circuits: Derandomizing the XOR Lemma.
\textit{Proceedings of the 29th ACM STOC}, pp.\ 220--229.
DOI: 10.1145/258533.258590.

\bibitem{FST-Unified2026}
L.~Geiger (2026).
The Renormalized Free Energy Framework.
Zenodo preprint.
\href{https://doi.org/10.5281/zenodo.19036190}{doi:10.5281/zenodo.19036190}.

\bibitem{GeigerRH2026c}
L.~Geiger (2026).
The Riemann Hypothesis Part III: The Analytical Rank-Finite Certificate.
Zenodo preprint.
\href{https://doi.org/10.5281/zenodo.19035640}{doi:10.5281/zenodo.19035640}.

\bibitem{Allender2006}
E.~Allender, H.~Buhrman, M.~Kouck\'y, D.~van~Melkebeek, \& D.~Ronneburger (2006).
Power from Random Strings.
\textit{SIAM Journal on Computing}, 35(6), 1467--1493.
DOI: 10.1137/050628994.

\bibitem{Hastad1987}
J.~H\r{a}stad (1987).
\textit{Computational Limitations of Small-Depth Circuits}.
MIT Press.

\bibitem{Cook1971}
S.~A.~Cook (1971).
The complexity of theorem-proving procedures.
\textit{Proceedings of the 3rd ACM Symposium on Theory of Computing}, pp.\ 151--158.
DOI: 10.1145/800157.805047.

\bibitem{ValiantVazirani1986}
L.~Valiant \& V.~Vazirani (1986).
NP is as easy as detecting unique solutions.
\textit{Theoretical Computer Science}, 47, 85--93.
DOI: 10.1016/0304-3975(86)90135-0.

\bibitem{Solomonoff1964}
R.~J.~Solomonoff (1964).
A formal theory of inductive inference, Parts I and II.
\textit{Information and Control}, 7(1), 1--22 and 7(2), 224--254.
DOI: 10.1016/S0019-9958(64)90131-7.

\bibitem{Allender2001}
E.~Allender (2001).
When worlds collide: Derandomization, lower bounds, and Kolmogorov complexity.
\textit{Proceedings of the 21st FSTTCS}, LNCS 2245, pp.\ 1--15.
DOI: 10.1007/3-540-45294-X\_1.

\bibitem{Hirahara2018}
S.~Hirahara (2018).
Non-black-box worst-case to average-case reductions within NP.
\textit{Proceedings of the 59th IEEE FOCS}, pp.\ 247--258.
DOI: 10.1109/FOCS.2018.00032.

\bibitem{LiuPass2020}
Y.~Liu and R.~Pass (2020).
On one-way functions and Kolmogorov complexity.
\textit{Proceedings of the 61st IEEE FOCS}, pp.\ 1243--1254.
DOI: 10.1109/FOCS46700.2020.00119. arXiv:2009.11514.

\bibitem{Impagliazzo1995}
R.~Impagliazzo (1995).
A personal view of average-case complexity.
\textit{Proceedings of the 10th Structure in Complexity Theory Conference}, pp.\ 134--147.
DOI: 10.1109/SCT.1995.514853.

\bibitem{AchlioptasRicci2006}
D.~Achlioptas, F.~Ricci-Tersenghi (2006).
On the solution-space geometry of random constraint satisfaction problems.
\textit{Proceedings of the 38th Annual ACM Symposium on Theory of Computing (STOC '06)}, pp.\ 130--139.
DOI: 10.1145/1132516.1132537.

\bibitem{MezardZecchina2002}
M.~M\'ezard, R.~Zecchina (2002).
Random $K$-satisfiability problem: From an analytic solution to an efficient algorithm.
\textit{Physical Review E}, 66, 056126.
DOI: 10.1103/PhysRevE.66.056126.

\bibitem{FortnowSanthanam2011}
L.~Fortnow \& R.~Santhanam (2011).
Infeasibility of instance compression and succinct PCPs for NP.
\textit{Journal of Computer and System Sciences}, 77(1), 91--106.
DOI: 10.1016/j.jcss.2010.06.007.

\bibitem{DellMelkebeek2014}
H.~Dell \& D.~van~Melkebeek (2014).
Satisfiability allows no nontrivial sparsification unless the polynomial-time hierarchy collapses.
\textit{Journal of the ACM}, 61(4), Article~23.
DOI: 10.1145/2629620.

\bibitem{Drucker2015}
A.~Drucker (2015).
New limits to classical and quantum instance compression.
\textit{SIAM Journal on Computing}, 44(5), 1443--1479.
DOI: 10.1137/130927115.

\bibitem{AllenderHirahara2019}
E.~Allender \& S.~Hirahara (2019).
New insights on the (non-)hardness of circuit minimization and related problems.
\textit{ACM Transactions on Computation Theory}, 11(4), 27:1--27:27.
DOI: 10.1145/3349616.

\bibitem{NatarajanWright2019}
A.~Natarajan \& J.~Wright (2019).
NEEXP $\subseteq$ MIP*.
\textit{Proceedings of the 60th IEEE FOCS}, pp.\ 510--528.
DOI: 10.1109/FOCS.2019.00039.

\bibitem{JiNatarajanVidickWrightYuen2021}
Z.~Ji, A.~Natarajan, T.~Vidick, J.~Wright, \& H.~Yuen (2021).
MIP* = RE.
\textit{Communications of the ACM}, 64(11), 131--138.
DOI: 10.1145/3485628.

\bibitem{Ilango2025}
S.~Hirahara and R.~Ilango,
\textit{NP-hardness of the Minimum Circuit Size Problem from Well-Studied Assumptions},
FOCS Proceedings, 2025.

\bibitem{GamarnikSudan2017}
D.~Gamarnik and M.~Sudan, \emph{Limits of local algorithms over sparse random graphs}, Ann.\ Probab.\ \textbf{45} (2017), 2353--2376.

\bibitem{Gamarnik2021}
D.~Gamarnik, \emph{The overlap gap property: a topological barrier to optimizing over random structures}, Proc.\ Natl.\ Acad.\ Sci.\ \textbf{118}(41) (2021), e2108492118.

\bibitem{AchlioptasCO2008}
D.~Achlioptas and A.~Coja-Oghlan, \emph{Algorithmic barriers from phase transitions}, Proc.\ 49th FOCS (2008), 793--802.

\bibitem{GamarnikJW2020}
D.~Gamarnik, A.~Jagannath, and A.~S.~Wein, \emph{The overlap gap property: a topological barrier to optimizing over random structures}, preprint, arXiv:2004.12063 (2020).

\bibitem{BreslerHuang2022}
G.~Bresler and B.~Huang, \emph{The algorithmic phase transition of random $k$-SAT for low degree polynomials},
\textit{Proceedings of the 63rd IEEE FOCS} (2022), pp.\ 298--309.
arXiv:2106.02129.

\bibitem{Fortnow2000}
L.~Fortnow (2000).
Kolmogorov complexity and computational complexity.
Vorlesungsnotizen, Mathematics Workshop, Kaikoura, Neuseeland; veröffentlicht als Überblick in \textit{Quaderni di Matematica}.

\bibitem{OrponenKoSchoeningWatanabe1994}
P.~Orponen, K.-I.~Ko, U.~Schöning, \& O.~Watanabe (1994).
Instance complexity.
\textit{Journal of the ACM}, 41(1), 96--121.
DOI: 10.1145/174644.174648.

\bibitem{Santhanam2023}
R.~Santhanam (2023).
An algorithmic approach to uniform lower bounds.
\textit{Proceedings of the 38th Computational Complexity Conference (CCC 2023)}, LIPIcs Vol.\ 264, Article~35.
DOI: 10.4230/LIPIcs.CCC.2023.35.

\bibitem[Geiger(2026d)]{Geiger2026-TU}
Geiger, L. (2026).
\newblock Functional Stability Theory: Turbulence and the K41 Spectrum as Free-Energy Minimum.
\newblock Preprint (FST-TU).

\bibitem[Geiger(2026e)]{Geiger2026-YM}
Geiger, L. (2026).
\newblock Functional Stability Theory: Volume-Independent Spectral Gap for Lattice Yang--Mills.
\newblock Preprint (FST-YM).

\end{thebibliography}

\end{document}
